\documentclass[11pt]{article}

\usepackage[letterpaper,margin=0.82in,headheight=15pt]{geometry}
\usepackage[T1]{fontenc}
\usepackage{lmodern}
\usepackage{microtype}
\usepackage{amsmath,amssymb,amsthm,bm,mathtools}
\usepackage{booktabs,longtable,tabularx,array}
\usepackage{enumitem}
\usepackage{xcolor}
\usepackage{fancyhdr}
\usepackage{graphicx}
\usepackage{tikz-cd}
\usepackage{hyperref}
\usepackage[nameinlink,noabbrev]{cleveref}

\definecolor{navy}{HTML}{17324D}
\definecolor{teal}{HTML}{147D92}
\definecolor{orange}{HTML}{D47A27}
\definecolor{softblue}{HTML}{EEF4F7}
\definecolor{softgray}{HTML}{F3F5F6}
\definecolor{darkgray}{HTML}{46545E}

\hypersetup{
  colorlinks=true,
  linkcolor=navy,
  citecolor=teal,
  urlcolor=teal,
  pdftitle={A rigorous algebraic and geometric architecture for a predictive spin-1 light-front model},
  pdfauthor={DeuteronWigner project}
}

\pagestyle{fancy}
\fancyhf{}
\fancyhead[L]{\small\color{darkgray}Predictive spin-1 light-front architecture}
\fancyhead[R]{\small\color{darkgray}Revised mathematical and physical specification}
\fancyfoot[L]{\small\color{darkgray}27 July 2026}
\fancyfoot[R]{\small\color{darkgray}\thepage}
\renewcommand{\headrulewidth}{0.4pt}

\setlist[itemize]{leftmargin=1.5em,itemsep=2pt,topsep=3pt}
\setlist[enumerate]{leftmargin=1.7em,itemsep=2pt,topsep=3pt}
\setlength{\parindent}{0pt}
\setlength{\parskip}{5pt}
\setlength{\emergencystretch}{2em}
\renewcommand{\arraystretch}{1.16}
\allowdisplaybreaks

\newtheorem{definition}{Definition}[section]
\newtheorem{principle}[definition]{Design principle}
\newtheorem{requirement}[definition]{Requirement}
\newtheorem{remark}[definition]{Remark}

\newcommand{\kT}{\bm{k}_{T}}
\newcommand{\pT}{\bm{p}_{T}}
\newcommand{\qT}{\bm{q}_{T}}
\newcommand{\DT}{\bm{\Delta}_{T}}
\newcommand{\DNT}{\bm{\Delta}_{N T}}
\newcommand{\bD}{\bm{b}_{\Delta}}
\newcommand{\bM}{\bm{b}_{\mathrm{TMD}}}
\newcommand{\RT}{\bm{R}_{T}}
\newcommand{\dd}{\mathrm{d}}
\newcommand{\Tr}{\operatorname{Tr}}
\newcommand{\Span}{\operatorname{span}}
\newcommand{\End}{\operatorname{End}}
\newcommand{\Hom}{\operatorname{Hom}}
\newcommand{\Obj}{\operatorname{Obj}}
\newcommand{\Id}{\operatorname{Id}}
\newcommand{\Hil}{\mathcal H}
\newcommand{\Kin}{\mathcal K}
\newcommand{\Alg}{\mathcal A}
\newcommand{\Rep}{\mathcal V}
\newcommand{\Path}{\mathsf{Path}}
\newcommand{\LF}{\mathrm{LF}}
\newcommand{\TMD}{\mathrm{TMD}}
\newcommand{\GTMD}{\mathrm{GTMD}}
\newcommand{\GPD}{\mathrm{GPD}}
\newcommand{\CP}{\mathrm{CP}}
\newcommand{\Amp}{\mathbf{Amp}}
\newcommand{\Dens}{\mathbf{Dens}}
\newcommand{\Match}{\mathbf{Match}}
\newcommand{\Red}{\mathbf{Red}}
\newcommand{\Proc}{\mathbf{Proc}}
\newcommand{\cO}{\mathcal O}
\newcommand{\cM}{\mathcal M}
\newcommand{\cR}{\mathcal R}
\newcommand{\cS}{\mathcal S}
\newcommand{\cV}{\mathcal V}
\newcommand{\cW}{\mathcal W}
\newcommand{\cE}{\mathcal E}
\newcommand{\cG}{\mathcal G}
\newcommand{\cQ}{\mathcal Q}
\newcommand{\cP}{\mathcal P}
\newcommand{\cC}{\mathcal C}
\newcommand{\cZ}{\mathcal Z}
\newcommand{\eps}{\varepsilon}

\newcolumntype{Y}{>{\raggedright\arraybackslash}X}
\newcolumntype{L}[1]{>{\raggedright\arraybackslash}p{#1}}

\title{\vspace{-1.0cm}
{\color{navy}\bfseries A Rigorous Algebraic and Geometric Architecture\\[2mm]
for a Predictive Spin-1 Light-Front Model}\\[4mm]
\large Typed amplitudes, gauge-covariant operators, matched nuclear composition,
and auditable reductions from microscopic dynamics to observables}
\author{DeuteronWigner project}
\date{27 July 2026}

\begin{document}
\maketitle

\begin{center}
\begin{minipage}{0.94\textwidth}
\colorbox{softblue}{\parbox{0.96\textwidth}{
\textbf{Status and purpose.}
This note is the corrected architecture specification for the model class that
would replace the present canonical phenomenological boundary by a genuinely
generative calculation.  It does not claim that algebra, geometry, category
theory, or topology supplies missing QCD dynamics.  It specifies how a
regulated light-front calculation, its gauge-link operators, its matched
nuclear reduction, and its observable maps must be typed so that exact
relations, approximations, and failures remain visible.  The primary
realizable route adopted here is a matched hierarchy from quark--gluon
nucleon dynamics to a spin-resolved nuclear light-front effective theory.
A direct quark--gluon deuteron solution remains an admissible later limit, but
is not assumed at the first implementation stage.
}}
\end{minipage}
\end{center}

\begin{abstract}
The current DeuteronWigner construction is a symmetry-complete and
nuclear-dressed phenomenological synthesis at a fixed boundary scale.  Its
common correlator contract preserves flavor, target polarization, transverse
rank, gauge-link class, nuclear mechanism, and uncertainty provenance, but
many scalar amplitudes are inherited from distinct fits or model sectors.
The next model class must instead generate many quark, antiquark, gluon,
nucleon, and deuteron observables from a restricted set of shared microscopic
interactions and matched operators.

We formulate the required calculation as a gauge-covariant, graded
light-front tensor architecture rather than as one undifferentiated
``category of everything.''  Regulated amplitude spaces and equivariant
linear maps form one layer; density operators and completely positive
reductions form another; regulator- and scheme-changing matching maps form a
third; GTMD, TMD, GPD, current, and process reductions form typed linear maps
between explicitly declared operator spaces.  Bare Wilson transport is a
functor from an admissible path groupoid to color-bundle fibers, whereas the
renormalized TMD operator additionally carries cusp, rapidity-regulator,
soft-subtraction, scale, and scheme data.  The transverse Wigner coordinate
\(\bD\), conjugate to \(\DT\), is kept strictly distinct from the TMD
impact variable \(\bM\), conjugate to \(\kT\).  Positivity is imposed on the
forward amplitude-level helicity correlator where a Gram representation is
valid, not on Wigner quasidistributions and not automatically on every
finite-order subtracted scheme object.

For the deuteron, the recommended first realization is a matched hierarchy:
a regulated quark--gluon light-front model generates nucleon amplitudes and
partonic operators; these are matched onto a spin-resolved nuclear
light-front effective theory with \(N\), \(\pi\), \(\Delta\), and any
retained non-nucleonic sectors; the full nuclear amplitude is evaluated
before any justified partial trace.  Tensor networks provide a sparse
symmetry-preserving contraction representation.  A cellular provenance
2-complex provides an audit of alternative paths, subtraction identities,
and double counting, but its homology is not asserted to determine physical
amplitudes.  A directed tower of regulated calculations and inverse matching
maps defines convergence at the level of renormalized observables.  The note
ends with software objects, property-based tests, staged milestones, and
acceptance conditions for calling the resulting model predictive.
\end{abstract}

\tableofcontents

\section{Scientific problem and change of model class}
\label{sec:problem}

\subsection{Present boundary and required direction of inference}

The present canonical model is scientifically useful because it assembles
physically informed inputs under a common correlator and composition
contract.  Schematically,
\begin{equation}
\{\text{PDF/TMD fits, wave functions, response models, phase models}\}
\longrightarrow
W^{\mathrm{assembled}}_{a/D}
\longrightarrow
\{F_A,\cO_{\mathrm{process}}\}.
\label{eq:assembled-direction}
\end{equation}
The resulting projections can be mutually consistent without being generated
by one microscopic state.  Their correlations are therefore partly
structural and partly inherited from the assumptions used to assemble the
inputs.

The next model must reverse the inference.  The primary implementation route
adopted in this note is
\begin{equation}
\begin{aligned}
&\bigl\{
H^{qg}_{\LF},\cR_{qg},\theta_{qg};
\cZ_{qg\to\mathrm{had}};
H^{\mathrm{nuc}}_{\LF},\cR_{\mathrm{nuc}},\theta_{\mathrm{nuc}};
\cZ_{\LF\to\mathrm{QCD}}
\bigr\}
\\
&\hspace{1.5cm}\longrightarrow
\bigl\{|\Psi_N\rangle,|\Psi_D\rangle,
\cO^{\mathrm{matched}}_{a/N},
\cO^{\mathrm{matched}}_{a/D},U_\gamma\bigr\}
\\
&\hspace{1.5cm}\longrightarrow
\bigl\{W_{a/N},W_{a/D}\bigr\}
\longrightarrow
\bigl\{F_A,\cO_{\mathrm{process}}\bigr\}.
\end{aligned}
\label{eq:matched-generative-direction}
\end{equation}
Here \(H^{qg}_{\LF}\) is a regulated quark--gluon light-front model for the
nucleon, \(H^{\mathrm{nuc}}_{\LF}\) is a consistently matched nuclear
light-front Hamiltonian or mass operator, and \(\cZ_{qg\to\mathrm{had}}\)
contains the operator and interaction matching required to pass from
quark--gluon degrees of freedom to a hadronic nuclear basis.  The map
\(\cZ_{\LF\to\mathrm{QCD}}\) matches the regulated Hamiltonian operators to
the declared renormalized QCD/GTMD/TMD scheme.

A direct calculation in which one quark--gluon Hamiltonian produces both the
nucleon and the six-quark deuteron is also admissible.  It is treated as a
future ultraviolet completion or cross-check of the matched hierarchy, not
as an assumption required for the first predictive implementation.

\subsection{What ``predictive'' means}

\begin{definition}[Predictive in this project]
A quantity is a prediction only when its normalization and functional form
are generated from shared calibrated interactions, states, currents, and
matching coefficients, and when the quantity or observable family was not
used directly to fit an independent coefficient with the same functional
freedom.
\end{definition}

This definition allows a controlled effective-QCD prediction without claiming
an exact solution of continuum QCD.  Predictive content comes from using
substantially fewer shared microscopic degrees of freedom than the number of
observables produced, from controlled regulator and truncation studies, and
from posterior predictive success on withheld channels.

\subsection{Design criteria}

\begin{principle}[Physical utility of mathematical structure]
A mathematical structure belongs in the production architecture only if it
performs at least one of the following tasks:
\begin{enumerate}
\item makes an exact symmetry, gauge constraint, support condition, or
operator identity automatic;
\item forbids an ill-typed composition, scheme mismatch, or double count;
\item embeds an approximation in a stated convergence or matching sequence;
\item forces testable relations among observables derived from a common state
or operator;
\item reduces redundant parameterization without deleting a physical degree
of freedom;
\item makes the source of a failed validation test identifiable.
\end{enumerate}
Terminology that performs none of these tasks is excluded.
\end{principle}

\begin{principle}[Architecture does not replace dynamics]
The architecture constrains where masses, couplings, counterterms, currents,
Wilson-line interactions, and nuclear operators may enter.  It does not
determine their values, prove factorization for an unsupported process, or
turn a finite matrix diagonalization into QCD without renormalization and
matching.
\end{principle}

\section{Convention and type contract}
\label{sec:conventions}

The convention layer is not descriptive metadata.  It is part of every state,
operator, transform, and matching-map type.  Code must fail closed when two
objects with incompatible conventions are composed.

\subsection{Light-front coordinates and the invariant mass operator}

Following the declared light-front convention used by the parent framework and standard front-form treatments \cite{Dirac1949,Brodsky1998,Vary2009}, we use
\begin{equation}
 g^{\mu\nu}=\operatorname{diag}(1,-1,-1,-1),
 \qquad
 v^{\pm}=\frac{v^0\pm v^3}{\sqrt2},
\label{eq:lf-components}
\end{equation}
so that
\begin{equation}
 v\cdot w=v^+w^-+v^-w^+-\bm v_T\cdot\bm w_T.
\label{eq:lf-dot-product}
\end{equation}
For null vectors \(n^2=\bar n^2=0\), \(n\cdot\bar n=1\),
\begin{equation}
 v^\mu=v^+\bar n^\mu+v^-n^\mu+v_T^\mu,
 \qquad
 v^2=2v^+v^- -\bm v_T^2.
\label{eq:lf-decomposition}
\end{equation}
The light-front time-translation generator is \(P^-\).  The invariant
mass-squared operator at fixed total momentum is
\begin{equation}
 \cM_{\LF}^2
 \equiv 2P^+P^- -\bm P_T^2.
\label{eq:lf-mass-operator}
\end{equation}
Accordingly, the eigenproblem is
\begin{equation}
 \cM_{\LF,\mathrm{ren}}^2
 |\Psi_h(P^+,\bm P_T,\Lambda)\rangle
 =M_h^2|\Psi_h(P^+,\bm P_T,\Lambda)\rangle.
\label{eq:lf-eigenproblem}
\end{equation}
Writing \(P^+P^-\) without the factor of two would correspond to a different
normalization of \(v^\pm\) and is disallowed unless the entire convention
contract is changed consistently.

\subsection{Distinct transverse variables}

At least three different transverse Fourier transforms and two different
Wigner phase-space pairs occur.  They must never share a software type:
\begin{longtable}{L{0.17\textwidth}L{0.18\textwidth}L{0.56\textwidth}}
\toprule
Variable & Fourier partner & Physical role\\
\midrule
\endhead
\(\DT\) & \(\bD\) & External hadron transverse momentum transfer and
partonic impact coordinate in the GTMD/Wigner transform at fixed skewness.
The partonic Wigner phase-space pair is \((\bD,\kT)\).\\
\(\kT\) & \(\bM\) & Average intrinsic parton transverse momentum and
bilocal transverse separation used in TMD Fourier transforms, small-\(b\)
matching, and Collins--Soper evolution.  \(\bM\) is not the Wigner impact
coordinate.\\
\(\DNT\) & \(\RT\) & Off-forward transverse transfer in the nuclear
spectral density and the corresponding active-nucleon impact coordinate.  The
nuclear Wigner phase-space pair is \((\RT,\pT)\), where \(\pT\) is the
average internal nucleon momentum retained during the \(\DNT\) transform.\\
\(\qT\) & process impact variable & Measured process-level transverse
momentum after the factorization convolution.  It is not identified with a
single parton's \(\kT\).\\
\bottomrule
\end{longtable}
The nuclear transfer \(\DNT\) is related to the external hadron transfer
\(\DT\) by the declared recoil map in the nuclear convolution.  It is stored
separately so that a factor such as \(1-y\), or another frame-dependent recoil
coefficient, cannot be inserted silently.

With the Fourier conventions
\begin{align}
 \rho_W(x,\kT,\bD)
 &=\int\frac{\dd^2\DT}{(2\pi)^2}
 e^{-i\bD\cdot\DT}
 W(x,\kT,\DT),
\label{eq:wigner-transform-convention}\\
 \widetilde F(x,\bM)
 &=\int\dd^2\kT\,e^{+i\bM\cdot\kT}F(x,\kT),
\label{eq:tmd-transform-convention}\\
 \rho_{N/D}(y,\pT,\RT)
 &=\int\frac{\dd^2\DNT}{(2\pi)^2}
 e^{-i\RT\cdot\DNT}
 \rho_{N/D}(y,\pT,\DNT),
\label{eq:nuclear-wigner-transform-convention}
\end{align}
\(\bD\) and \(\RT\) are Wigner-space coordinates in different physical
spaces, while \(\bM\) is the transverse separation of the bilocal TMD
operator.  Equal dimensions do not imply equal meaning.

\subsection{Operator identity is decorated, not nominal}

A production operator is not identified only by a name such as
``quark GTMD'' or ``Sivers operator.''  Its type contains at least
\begin{equation}
 \mathfrak o=
 \bigl(
 a,\Gamma,\xi,\Delta,\gamma,R_{\mathrm{color}},
 \cR_{\mathrm{UV}},\cR_{\mathrm{rap}},
 \cS_{\mathrm{soft}},\mu,\zeta,\mathfrak s
 \bigr),
\label{eq:operator-type-tuple}
\end{equation}
where \(a\) is the parton species, \(\Gamma\) is the quark Dirac or gluon
Lorentz projection, \(\gamma\) is the Wilson path, \(R_{\mathrm{color}}\)
is its color representation, \(\cR_{\mathrm{UV}}\) and
\(\cR_{\mathrm{rap}}\) are regulators, \(\cS_{\mathrm{soft}}\) specifies
soft subtraction, and \(\mathfrak s\) is the renormalization/factorization
scheme.  Cusp geometry, transverse closure, staple orientation, and the pair
of gluon links are included in \(\gamma\).

\begin{requirement}[Fail-closed composition]
Two operators, states, or maps may compose only if their momentum convention,
normalization, color representation, gauge-link class, regulator, subtraction
scheme, and domain/codomain labels agree or if an explicit matching adapter
is supplied.
\end{requirement}

\subsection{Map classes must remain distinct}

The calculation uses several composition laws.  They are related, but they
are not all morphisms of one undifferentiated category:
\begin{longtable}{L{0.16\textwidth}L{0.28\textwidth}L{0.46\textwidth}}
\toprule
Layer & Typical objects and maps & Meaning\\
\midrule
\endhead
\(\Amp\) & Graded regulated Hilbert spaces; equivariant linear amplitude
maps; adjoints & Hamiltonian evolution, interaction vertices, current
insertions, constituent coupling, coherent amplitudes.\\
\(\Dens\) & Matrix/operator algebras with density matrices as states; positive or completely
positive trace-nonincreasing maps & Partial traces, detector/environment
reduction, incoherent channel composition after amplitudes are formed.\\
\(\Match\) & Regulator- and scheme-indexed operator spaces; matching and
coarse-graining maps & Renormalization, EFT matching, soft subtraction,
small-\(b\) OPE, basis/truncation comparison.\\
\(\Red\) & GTMD, TMD, GPD, current, and Wigner spaces; typed linear limits,
integrals, Fourier transforms, and moments & Observable and marginal
reductions of a common matched operator.\\
\(\Proc\) & Factorized hadronic tensors and cross sections; hard, soft,
fragmentation, and fixed-order kernels & Process-specific assembly, including
factorization-status and \(W+Y\) information.\\
\bottomrule
\end{longtable}
At finite truncation, \(\Amp\) can be implemented as a dagger symmetric
monoidal category of block-sparse finite-dimensional spaces and equivariant
linear maps.  The tensor product is first a kinematical constituent product;
gauge and total-quantum-number projectors then select the physical subspace.
Continuum unbounded-operator domain questions are handled through the
truncation tower rather than hidden inside categorical notation.

\section{Regulated light-front state spaces and dynamics}
\label{sec:hilbert}

\subsection{Kinematical sectors and physical subspaces}

For a retained Fock label
\begin{equation}
 \nu=(n_q,n_{\bar q},n_g;B,Q,S,\ldots),
\end{equation}
the kinematical sector is
\begin{equation}
 \Kin_\nu=
 L^2(\cM_\nu,\dd\mu_\nu)
 \otimes
 \bigotimes_{i\in\nu}
 \left(
 \Rep_{f_i}\otimes\Rep_{c_i}\otimes\Rep_{\lambda_i}
 \right)
 \otimes\Rep_{L_z},
\label{eq:kinematical-sector}
\end{equation}
with
\begin{equation}
 \cM_\nu=
 \left\{(x_i,\bm k_{iT})\;\middle|\;
 x_i>0,\quad
 \sum_i x_i=1,\quad
 \sum_i\bm k_{iT}=0
 \right\}.
\label{eq:momentum-simplex}
\end{equation}
The physical sector is not the full tensor product.  It is the simultaneous
solution space of the gauge, statistics, support, and total-quantum-number
constraints,
\begin{equation}
 \Hil_\nu=
 \left\{|\psi\rangle\in\Kin_\nu\;\middle|\;
 \cG^a|\psi\rangle=0,\quad
 \cQ_i|\psi\rangle=q_i|\psi\rangle,\quad
 \Pi_{\mathrm{statistics}}|\psi\rangle=|\psi\rangle,
 \ \ldots\right\}
 =\Pi_{\mathrm{phys},\nu}\Kin_\nu.
\label{eq:physical-sector}
\end{equation}
Here \(\cG^a\) denotes the regulated Gauss constraint; for an isolated hadron,
the selected physical sector is globally color singlet.  The notation
\(\Pi_{\mathrm{phys},\nu}\) does not assume that separately implemented
constraint projectors commute before their joint solution is constructed.
The regulated state space is the direct sum
\begin{equation}
 \Hil_{\LF}^{(r)}=
 \bigoplus_{\nu\in\mathcal F_r}\Hil_\nu^{(r)},
\label{eq:regulated-direct-sum}
\end{equation}
where \(r\) records basis, volume, ultraviolet, infrared, longitudinal, and
Fock-sector truncations.

For a nucleon pilot, the minimum useful sector content is
\begin{equation}
 |N,\lambda\rangle=
 \psi_{qqq}^{\lambda}|qqq\rangle
 +\psi_{qqqg}^{\lambda}|qqqg\rangle
 +\psi_{qqqq\bar q}^{\lambda}|qqqq\bar q\rangle
 +\cdots.
\label{eq:nucleon-fock-expansion}
\end{equation}
A valence-only sector cannot generate antiquark distributions, a dynamical
gluon correlator, genuine quark--gluon matrix elements, or the complete
Wilson-line rescattering structure.

\subsection{Renormalized block mass operator}

Let \(P_\nu\) be the projector onto sector \(\Hil_\nu^{(r)}\).  The regulated
mass-squared operator has blocks
\begin{equation}
 \cM_{\nu\mu}^{2(r)}
 =P_\nu\cM_{\LF}^{2(r)}P_\mu.
\label{eq:mass-blocks}
\end{equation}
A schematic nucleon block structure is
\begin{equation}
\cM_{\LF}^{2(r)}\sim
\begin{pmatrix}
 \cM^2_{qqq} & V_{qqq\leftrightarrow qqqg}
 &V_{qqq\leftrightarrow qqqq\bar q}&\cdots\\
 V^\dagger_{qqq\leftrightarrow qqqg}&\cM^2_{qqqg}
 &V_{qqqg\leftrightarrow qqqq\bar q}&\cdots\\
 V^\dagger_{qqq\leftrightarrow qqqq\bar q}
 &V^\dagger_{qqqg\leftrightarrow qqqq\bar q}
 &\cM^2_{qqqq\bar q}&\cdots\\
 \vdots&\vdots&\vdots&\ddots
\end{pmatrix}.
\label{eq:block-mass-operator}
\end{equation}
The diagonal blocks contain kinetic terms and sector-preserving interactions.
Off-diagonal blocks generate emission, absorption, and pair creation.  The
operator definition must declare quark masses, confining or effective
interactions, quark--gluon vertices, instantaneous light-front terms,
zero-mode treatment, counterterms, and renormalization conditions.

\begin{requirement}[Renormalization is part of the model]
A finite matrix and a smooth eigenvector are not by themselves a predictive
Hamiltonian model.  The regulator, counterterm basis, renormalization
conditions, omitted-sector effective operators, and convergence sequence are
part of the model definition and versioned output.
\end{requirement}

\subsection{Light-front covariance and rotational restoration}

The front form separates kinematical and dynamical Poincar\'e generators \cite{Dirac1949,Brodsky1998}.
The longitudinal boost structure and \(J^z\) labeling are especially useful,
but full rotational covariance is not guaranteed by attaching an
\(SU(2)\) spin label to a truncated basis.  The calculation must test the
Poincar\'e commutator defects accessible in the chosen representation,
rotational multiplet splittings, current angular conditions, and their
behavior under basis and Fock-sector enlargement.

Fock probabilities, separate quark/gluon spin decompositions, and canonical
OAM components can be regulator, gauge, and scheme dependent.  They are
reported as state diagnostics within a declared scheme; convergence claims
of physical observables must not be replaced by apparent stability of these
non-observable decompositions.

\section{Gauge, internal symmetry, and representation grammar}
\label{sec:symmetry}

\subsection{Do not package unlike structures into one exact direct product}

The production architecture distinguishes:
\begin{enumerate}
\item the local \(SU(3)_{\mathrm{color}}\) gauge action and Gauss-law
constraint;
\item exact conserved labels retained by the regulated model, such as baryon
number and electric charge;
\item approximate flavor/isospin symmetry and its explicit breaking by
\(m_d-m_u\), QED, and nuclear interactions;
\item the light-front kinematical subgroup and the dynamical restoration of
full Poincar\'e covariance;
\item discrete parity, time reversal, and charge conjugation, including their
action on Wilson paths and external states.
\end{enumerate}
It is therefore misleading to write one exact global group
\(SU(3)_{\mathrm{color}}\times SU(2)_{I}\times SU(2)_{\mathrm{spin}}
\times SO(2)_{L_z}\) and treat every factor identically.

\subsection{Intertwiners define allowed couplings}

For a declared symmetry action, a coupling tensor is an equivariant map
\begin{equation}
 C:\Rep_{R_1}\otimes\cdots\otimes\Rep_{R_n}
 \longrightarrow\Rep_{R_{\mathrm{out}}},
 \qquad
 C D_{\mathrm{in}}(g)=D_{\mathrm{out}}(g)C.
\label{eq:intertwiner-definition}
\end{equation}
Only invariant or correctly covariant tensors are parameterized.  Color
singlet couplings, flavor tensors, spin recouplings, and transverse angular
weights are therefore restrictions on the allowed map space rather than
penalties applied after a fit.

For gluon T-odd sectors, the independent invariant tensors
\begin{equation}
 f^{abc},\qquad d^{abc}
\label{eq:f-d-tensors}
\end{equation}
are different intertwiner channels.  Their amplitudes cannot be identified by
a universal scalar coefficient, and a process kernel must specify the linear
combination it probes.

\subsection{Spin-1 density irreducibles}

For the spin-1 target representation \(\Rep_1\), consistent with the standard spin-1 correlator decompositions \cite{BacchettaMulders2000,KumanoSong2021,Boer2016},
\begin{equation}
 \End(\Rep_1)
 \cong \Rep_1\otimes\Rep_1^*
 \cong \Rep_0\oplus\Rep_1\oplus\Rep_2.
\label{eq:spin1-end-decomposition}
\end{equation}
Thus a target density operator has scalar, vector, and quadrupole pieces,
\begin{equation}
 \rho_D=
 \rho_D^{(0)}
 +\sum_{m=-1}^{1}S_m T_m^{(1)}
 +\sum_{m=-2}^{2}Q_m T_m^{(2)}.
\label{eq:spin1-density-decomposition}
\end{equation}
The familiar \(U,L,T,LL,LT,TT\) components are fixed real Cartesian
combinations of these spherical tensors.  A convention adapter maps between
the irreducible basis and any Trento-like or project-specific normalization;
the sign of \(LL\) is never inferred from a function name.

This representation makes rotations exact at the algebraic level,
separates vector from tensor response, and gives pure-\(S\), \(S\)--\(D\),
and higher-OAM limits explicit Clebsch--Gordan tests.

\subsection{Transverse rank and angular weight}

The transverse rotation group \(SO(2)\) has one-dimensional irreducible
weights \(m\in\mathbb Z\).  A rank-\(|m|\) harmonic carries an azimuthal
factor \(e^{im\phi}\) or its real symmetric-traceless equivalent.  A matrix
element is allowed only when the total light-front \(J^z\) weight of the
initial state, operator tensor, parton-polarization block, and final state
closes.  The sign of the integer relation depends on the helicity-state phase
and on whether the parton-helicity change is included in the operator label;
therefore the implementation derives the relation from equivariance rather
than hard-coding a convention-independent sign formula.

The practical zero test is unambiguous: if a named projection requires an
interference between two amplitude blocks with different \(L_z\), helicity,
or target-spin labels, disabling either block must remove the corresponding
angular harmonic.

\subsection{Rank-aware Fourier--Bessel transforms}

A transverse harmonic of weight \(m\) transforms radially with
\(J_{|m|}\).  With convention-dependent mass, phase, and tensor factors
stored in \(\mathcal N_m\) and \(\widetilde{\mathcal N}_m\),
\begin{align}
 \widetilde F^{(m)}(x,b_{\mathrm{TMD}})
 &=2\pi\int_0^\infty k_T\,\dd k_T\,
 J_{|m|}(b_{\mathrm{TMD}}k_T)
 \mathcal N_m(k_T)F^{(m)}(x,k_T),
\label{eq:rank-bessel-forward}\\
 F^{(m)}(x,k_T)
 &=\int_0^\infty\frac{b_{\mathrm{TMD}}\,\dd b_{\mathrm{TMD}}}{2\pi}
 J_{|m|}(b_{\mathrm{TMD}}k_T)
 \widetilde{\mathcal N}_m(b_{\mathrm{TMD}})
 \widetilde F^{(m)}(x,b_{\mathrm{TMD}}).
\label{eq:rank-bessel-inverse}
\end{align}
The registry stores \(m\), the tensor basis, the extracted powers of
\(k_T/M\) or \(b_{\mathrm{TMD}}M\), and any Fourier phase.  A scalar
\(J_0\) transform may not be applied to rank-one or rank-two amplitudes.

\subsection{Controlled isospin}

A useful decomposition is
\begin{equation}
 \cM_{\LF}^2
 =\cM_0^2
 +\delta\cM_{m_d-m_u}^2
 +\delta\cM_{\mathrm{QED}}^2
 +\delta\cM_{\mathrm{nuc}}^2.
\label{eq:isospin-breaking-mass}
\end{equation}
Exact isospin is recovered only when the breaking operators are disabled.
It is not implemented by setting proton \(u\) and \(d\) amplitudes equal.
Proton and neutron remain distinct states in an isospin multiplet, allowing
both an exact charge-symmetry test and realistic charge-symmetry breaking.

\section{Typed categorical architecture}
\label{sec:typed-categories}

\subsection{Amplitude layer}

At a finite truncation, define \(\Amp_r\) as follows, using ordinary categorical composition only where domains and codomains are genuinely compatible \cite{MacLane1998,Selinger2007}:
\begin{itemize}
\item objects are regulated graded physical Hilbert spaces with momentum,
normalization, gauge, and symmetry metadata;
\item morphisms are admissible equivariant linear maps with declared domains
and codomains;
\item the dagger is the Hilbert-space adjoint;
\item the monoidal product is a kinematical constituent product followed by
explicit physical fusion/projector maps;
\item direct sums encode alternative Fock sectors.
\end{itemize}
Hamiltonian vertices, current insertions, Wilson-line interaction vertices,
and nuclear sector couplers live in this layer.

\subsection{Density and channel layer}

At finite truncation, objects of \(\Dens_r\) are matrix algebras attached to the relevant Hilbert spaces; density operators are positive trace-class states in those algebras.  A map
\(\cE:\Alg_A\to\Alg_B\) is completely positive when
\(\cE\otimes\Id_n\) is positive for every auxiliary matrix dimension
\(n\).  Trace-nonincreasing CP maps describe a selected incoherent channel;
trace-preserving maps require a complete set of retained outcomes.

Complete positivity is not a property of an amplitude map.  It arises after
an amplitude embedding and a justified partial trace, as described in
\cref{sec:nuclear-composition}.

\subsection{Matching layer}

\(\Match\) contains regulator-, basis-, resolution-, and scheme-indexed
operator spaces.  A matching morphism
\begin{equation}
 \cZ_{A\to B}:\cO_A\longrightarrow\cO_B
\label{eq:generic-matching-map}
\end{equation}
is not physical time evolution.  It changes the description while preserving
a declared set of renormalized matrix elements up to quantified truncation
errors.  Examples include quark--gluon-to-hadronic EFT matching, omitted-Fock-
sector counterterms, UV renormalization, rapidity renormalization, soft
subtraction, and small-\(\bM\) operator-product matching.

\subsection{Reduction and process layers}

\(\Red\) contains linear operations such as a forward limit, a Fourier
transform, a transverse-momentum integration, or a local moment.  Such a map
is defined only on an operator space for which the corresponding limit is
renormalized and meaningful.  \(\Proc\) then combines matched distributions
with hard, fragmentation, soft, color, and fixed-order kernels for a specified
process and factorization status.

\subsection{Commuting squares rather than one universal composition law}

The essential consistency statements are mixed squares.  For example,
\begin{equation}
\begin{tikzcd}[column sep=large,row sep=large]
 W_{\LF}^{(r)}
 \arrow[r,"\cZ_{\LF\to\mathrm{QCD}}"]
 \arrow[d,"\mathsf R_{\LF}"']
 & W_{\mathrm{QCD}}^{\mathfrak s}(\mu,\zeta)
 \arrow[d,"\mathsf R_{\mathrm{QCD}}^{\mathfrak s}"]\\
 R_{\LF}^{(r)}
 \arrow[r,"\cZ_{\LF\to R}"']
 & R_{\mathrm{QCD}}^{\mathfrak s}(\mu,\zeta)
\end{tikzcd}
\label{eq:matched-reduction-square}
\end{equation}
commutes only within the declared perturbative, regulator, basis, and
numerical error,
\begin{equation}
 \left\|
 \mathsf R_{\mathrm{QCD}}^{\mathfrak s}
 \cZ_{\LF\to\mathrm{QCD}}W_{\LF}^{(r)}
 -
 \cZ_{\LF\to R}\mathsf R_{\LF}W_{\LF}^{(r)}
 \right\|
 \leq\eps_{\mathrm{match}}+\eps_{\mathrm{trunc}}+\eps_{\mathrm{num}}.
\label{eq:commuting-error-bound}
\end{equation}
The square is a validation condition, not a formal equality that erases
matching coefficients.

\section{Gauge links as bare transport and renormalized operators}
\label{sec:gauge-links}

\subsection{Connection, fibers, and admissible paths}

Let \(P(M,SU(3))\) be the color principal bundle over the spacetime domain
used by the regulated operator, with connection
\begin{equation}
 A=A_\mu^a t^a\,\dd x^\mu.
\end{equation}
For a representation \(R\), the associated vector bundle is
\(E_R=P\times_R\Rep_R\).  If \(\gamma:a\to b\) is an oriented admissible
path, bare parallel transport is
\begin{equation}
 U_R(\gamma):E_{R,a}\longrightarrow E_{R,b},
 \qquad
 U_R(\gamma)
 =\mathcal P\exp\left[-ig\int_\gamma A_R\right].
\label{eq:bare-holonomy}
\end{equation}
It transforms gauge-covariantly,
\begin{equation}
 U_R(\gamma)\longmapsto
 G_R(b)U_R(\gamma)G_R(a)^{-1}.
\label{eq:holonomy-gauge-transform}
\end{equation}
Thus an open Wilson line is a map between endpoint fibers, not an
gauge-invariant endomorphism by itself.

The appropriate geometric object is the path groupoid, or more precisely a
path groupoid modulo reparametrization and any declared thin-homotopy
identification:
\begin{align}
 \Obj(\Path)&=\{\text{operator endpoints}\},\\
 \Hom_{\Path}(a,b)&=\{\text{admissible oriented paths }\gamma:a\to b\}.
\end{align}
Bare transport is a functor to color-bundle fiber isomorphisms,
\begin{equation}
 \mathcal U_R:\Path\longrightarrow\mathbf{VectIso},
 \qquad
 \gamma\longmapsto U_R(\gamma),
\label{eq:path-functor}
\end{equation}
with
\begin{equation}
 U_R(\gamma_2\circ\gamma_1)
 =U_R(\gamma_2)U_R(\gamma_1),
 \qquad
 U_R(\gamma^{-1})=U_R(\gamma)^{-1}.
\label{eq:bare-path-composition}
\end{equation}
For a unitary representation and a real gauge connection,
\(U_R(\gamma)^{-1}=U_R(\gamma)^\dagger\).

\subsection{Curvature and the limit of topological language}

The curvature
\begin{equation}
 F=\dd A-igA\wedge A
\label{eq:gauge-curvature}
\end{equation}
controls local path dependence.  In generic QCD configurations
\(F_{\mu\nu}\neq0\), so Wilson transport depends on path geometry and not
only on ordinary homotopy class.  The path groupoid records endpoints,
orientation, and allowed concatenation.  It does not make a Sivers,
Boer--Mulders, or gluon T-odd distribution a topological invariant.

\subsection{Staples, link reversal, and gluon link pairs}

A TMD staple contains longitudinal or near-lightlike segments, a transverse
closure, a direction toward future or past infinity, and a rapidity
regularization prescription.  Future- and past-pointing staples are distinct
paths.  They are not generally literal inverses of one another.  Time
reversal, together with the required complex conjugation and state
transformation, maps the corresponding operator classes and gives the
well-known process reversal relations when TMD factorization applies,
\begin{align}
 F_{\mathrm{even}}^{[+]}&=F_{\mathrm{even}}^{[-]},\\
 F_{\mathrm{odd}}^{[+]}&=-F_{\mathrm{odd}}^{[-]}.
\label{eq:t-even-odd-reversal}
\end{align}

Quark links act in the fundamental representation.  Gluon TMD operators
contain a pair of adjoint link structures, customarily labeled by classes
such as \([+,+]\), \([-,-]\), \([+,-]\), and \([-,+]\).  The independent
\(f^{abc}\)- and \(d^{abc}\)-type color contractions remain explicit in the
operator and process types.

\subsection{Bare transport is not the complete renormalized TMD operator}

Lightlike or near-lightlike Wilson lines have ultraviolet, cusp, and rapidity
divergences \cite{Collins2011,Echevarria2012}.  The physical operator therefore requires a decorated identity
\begin{equation}
 \mathfrak W=
 \bigl(
 \gamma,R,\text{cusp data},\cR_{\mathrm{UV}},\cR_{\mathrm{rap}},
 \cS_{\mathrm{soft}},\mu,\zeta,\mathfrak s
 \bigr).
\label{eq:decorated-wilson-operator}
\end{equation}
A schematic renormalized quark operator is
\begin{equation}
 \cO_{q}^{\mathfrak s}(\mu,\zeta)
 =Z_q^{\mathfrak s}(\mu,\zeta;\cR_{\mathrm{UV}},\cR_{\mathrm{rap}})
 \left[\cS_q^{\mathfrak s}\right]^{-1/2}
 \cO_{q}^{\mathrm{bare}}[U_R(\gamma)].
\label{eq:renormalized-tmd-operator}
\end{equation}
The bare path law in \cref{eq:bare-path-composition} remains exact for
transport.  Composition of renormalized cusped operators additionally
requires the declared renormalization and soft factors.  The software must
therefore keep \texttt{BareHolonomy} distinct from
\texttt{RenormalizedLinkOperator}.

\subsection{Where naive-T-odd dynamics enters}

Expanding the bare transport,
\begin{equation}
 U_R(\gamma)=1-ig\int_\gamma A_R
 -g^2\int_{\gamma}\!\int_{\gamma}
 \mathcal P[A_R(\xi)A_R(\eta)]+\cdots,
\label{eq:holonomy-expansion}
\end{equation}
shows where the gauge field enters, but a nonzero naive-T-odd distribution
requires absorptive interference: intermediate states, pole prescriptions,
rescattering, and a factorization-appropriate eikonal limit.  The path type
specifies which phase transformation is allowed.  It does not determine the
phase numerically.

The first dynamical Wilson-line pilot must demonstrate:
\begin{enumerate}
\item vanishing of naive-T-odd functions when rescattering is disabled;
\item future/past reversal in processes for which the theorem applies;
\item flavor and operator dependence from shared intermediate states;
\item separate gluon \(f\)- and \(d\)-type contractions;
\item a perturbative one-gluon limit and controlled higher-eikonal-order
variation;
\item explicit labels for established, conditional, exploratory, or broken
factorization in each proposed process.
\end{enumerate}

\section{GTMD operators and matched commuting reductions}
\label{sec:gtmd-reductions}

\subsection{Off-forward operator fibers}

A GTMD insertion changes the external hadron momentum and is the common off-forward parent of the TMD, GPD, and Wigner reductions \cite{Meissner2009,Belitsky2004,LorcePasquini2011}.  It is therefore most
precise to write an operator fiber
\begin{equation}
 \cW_{\Gamma,\mathfrak W}(P',P):
 \Hil_{h,P}\longrightarrow\Hil_{h,P'},
\label{eq:gtmd-operator-fiber}
\end{equation}
or equivalently to regard it as a sesquilinear form on a direct integral of
momentum fibers.  Its matrix element is
\begin{equation}
 W_{\Lambda'\Lambda}^{[\Gamma,\mathfrak W]}
 (x,\xi,\kT,\DT)
 =\langle P',\Lambda'|
 \cW_{\Gamma,\mathfrak W}(x,\xi,\kT,\DT)
 |P,\Lambda\rangle.
\label{eq:gtmd-matrix-element}
\end{equation}
For the first implementation, \(\xi=0\),
\(P'=P+\Delta/2\), \(P=P-\Delta/2\), and \(\Delta^+=0\).
Quark, antiquark, and gluon operators remain distinct, with their own
Dirac/Lorentz projections and color-link structures.

A light-front wave-function calculation evaluates
\cref{eq:gtmd-matrix-element} as sums of diagonal and, when the operator
changes resolved particle content, off-diagonal Fock-sector overlaps.  The
same state amplitudes must feed every compatible quark, antiquark, and gluon
projection.

\subsection{Hamiltonian-to-QCD matching}

A low-resolution overlap is not yet a TMD or GTMD in a named QCD scheme.  The
matching problem is schematically
\begin{equation}
 W_{a/h}^{\mathrm{QCD},\mathfrak s}
 (x,\xi,\bM,\DT;\mu,\zeta)
 =\sum_b
 \cZ_{a\leftarrow b}^{\LF\to\mathrm{QCD},\mathfrak s}
 \otimes
 W_{b/h}^{\LF,(r)}
 +\delta W_{a/h}^{(r)},
\label{eq:lf-to-qcd-gtmd-matching}
\end{equation}
where the matching kernel carries the Hamiltonian regulator, UV and rapidity
renormalization, soft subtraction, operator mixing, and scale dependence.
The remainder \(\delta W^{(r)}\) contains power, Fock-sector, basis, and
matching truncations and must be estimated separately.

\subsection{Four physically different reductions}

The following maps are distinct:
\begin{align}
 \mathsf T[W]
 &\equiv W\big|_{\xi=0,\DT=0},
 &&\text{forward TMD operator class},
\label{eq:tmd-reduction}\\
 \mathsf G[W]
 &\equiv \left[\int\dd^2\kT\,W\right]_{\mathrm{ren/matched}},
 &&\text{GPD operator class},
\label{eq:gpd-reduction}\\
 \mathsf W_{\Delta}[W]
 &\equiv\int\frac{\dd^2\DT}{(2\pi)^2}
 e^{-i\bD\cdot\DT}W,
 &&\text{partonic Wigner transform},
\label{eq:wigner-reduction}\\
 \mathsf M_n[W]
 &\equiv\left[\int\dd x\,x^{n-1}\int\dd^2\kT\,W\right]_{\mathrm{ren}},
 &&\text{local-current or EMT moment}.
\label{eq:local-moment-reduction}
\end{align}
The power of \(x\), overall normalization, and operator mixing in \cref{eq:local-moment-reduction} are species and operator dependent; the displayed expression is schematic.  The brackets in \cref{eq:gpd-reduction,eq:local-moment-reduction} are
essential.  A literal integral of a soft-subtracted TMD over all \(\kT\), or
a literal evaluation at \(\bM=0\), is ultraviolet singular in general.
The relation to a collinear operator is supplied by renormalized operator
matching or the small-\(\bM\) OPE, not by deleting the transverse separation
without qualification.

At small \(b_{\mathrm{TMD}}\),
\begin{equation}
 \widetilde F_{a/h}^{\mathfrak s}
 (x,\bM;\mu,\zeta)
 =\sum_b C_{a\leftarrow b}^{\mathfrak s}
 (\bM;\mu,\zeta)\otimes f_{b/h}(x;\mu)
 +\cO(b_{\mathrm{TMD}}\Lambda_{\mathrm{QCD}}),
\label{eq:small-b-ope}
\end{equation}
with operator- and transverse-rank-dependent coefficients.  This equation is
the correct bridge between a TMD and its collinear limit.

\subsection{Reduction closure is a matched statement}

The desired common-parent relation is not the naive equality
\(\int\dd^2\kT\,\mathsf T[W]=\mathsf M_1[W]\).  It is the statement that
all routes through the regulator, matching, and reduction diagram agree
within their declared errors.  For example,
\begin{equation}
 \mathsf M_1^{\mathrm{QCD}}
 \cZ_{\LF\to\mathrm{QCD}}[W_{\LF}^{(r)}]
 \simeq
 \cZ_{\LF\to J}
 \mathsf M_1^{\LF}[W_{\LF}^{(r)}],
\label{eq:current-closure-matched}
\end{equation}
and
\begin{equation}
 \mathsf T^{\mathfrak s}
 \cZ_{\LF\to\mathrm{QCD}}[W_{\LF}^{(r)}]
 \simeq
 \cZ_{\LF\to\TMD}^{\mathfrak s}
 \mathsf T^{\LF}[W_{\LF}^{(r)}].
\label{eq:tmd-closure-matched}
\end{equation}
The equality sign is recovered only at the specified order and in the
continuum/truncation limit.

\subsection{Nuclear composition and projection need not commute automatically}

Let \(\mathsf N\) denote the full spin-resolved nuclear amplitude map and
\(\mathsf P_A\) a named partonic projection.  The diagnostic square
\begin{equation}
 \mathsf P_A\circ\mathsf N[W_N]
 \stackrel{?}{\simeq}
 \mathsf N_A\circ\mathsf P[W_N]
\label{eq:nuclear-projection-square}
\end{equation}
can fail because a projection-dependent two-body current, coherent phase, or
operator matching term is absent.  Such a failure is not repaired by forcing
the square to commute algebraically.  It identifies the missing amplitude or
operator in the nuclear theory.

\section{Transverse phase-space geometry and orbital structure}
\label{sec:symplectic}

\subsection{The partonic Wigner phase-space variables}

At zero skewness, the transverse Wigner variables are
\begin{equation}
 (\bD,\kT)\in T^*\mathbb R^2.
\label{eq:wigner-phase-space}
\end{equation}
As a kinematic organizing manifold, this cotangent bundle carries the
canonical symplectic form
\begin{equation}
 \omega_\perp
 =\dd b_{\Delta x}\wedge\dd k_x
 +\dd b_{\Delta y}\wedge\dd k_y.
\label{eq:transverse-symplectic-form}
\end{equation}
The \(SO(2)\) rotation action has moment map
\begin{equation}
 \mu_{L_z}(\bD,\kT)
 =b_{\Delta x}k_y-b_{\Delta y}k_x.
\label{eq:transverse-moment-map}
\end{equation}
A Wigner-space OAM moment therefore has the schematic form
\begin{equation}
 \langle L_z\rangle_{\mathfrak W}
 =\int\dd x\,\dd^2\bD\,\dd^2\kT\,
 \mu_{L_z}(\bD,\kT)
 \rho_W(x,\bD,\kT;\mathfrak W).
\label{eq:wigner-oam-moment}
\end{equation}
The Wilson-path and operator labels are retained because canonical and
kinetic OAM correspond to different gauge-link/operator constructions \cite{Lorce2012}.

\begin{remark}[Scope of the symplectic statement]
\Cref{eq:transverse-symplectic-form} organizes the transverse Fourier and
rotation structure.  It is not, by itself, a derivation of the full reduced
symplectic form of interacting gauge theory after constraints and zero modes
are resolved.  The physical OAM interpretation comes from the matched QCD
operator, not from the notation \(T^*\mathbb R^2\) alone.
\end{remark}

\subsection{The TMD impact variable is not a phase-space position}

The pair \((\bM,\kT)\) is a Fourier pair associated with the bilocal operator
separation and its momentum representation.  It is not the Wigner
phase-space pair.  In particular,
\begin{equation}
 \mu_{L_z}(\bM,\kT)
\end{equation}
is not used as the partonic spatial OAM moment.  This distinction is enforced
by separate classes \texttt{WignerImpactCoordinate} and
\texttt{TMDSeparation}.

\subsection{Nuclear and partonic Wigner objects remain distinct}

The nuclear Wigner/spectral density
\begin{equation}
 \rho_{N/D}(y,\pT,\RT)
\label{eq:nuclear-wigner-density}
\end{equation}
is obtained by Fourier transformation in the nuclear off-forward transfer \(\DNT\) while retaining the average internal momentum \(\pT\), as in \cref{eq:nuclear-wigner-transform-convention}.  It describes nucleonic degrees of freedom in the deuteron and follows the standard phase-space logic of nuclear Wigner constructions \cite{Neff2016}.  The partonic Wigner
distribution
\begin{equation}
 \rho_{a/D}(x,\kT,\bD;\mathfrak W)
\label{eq:partonic-wigner-density}
\end{equation}
is a QCD operator transform.  They are composed by a nuclear amplitude or
matched operator map; they are not two names for the same object.

\subsection{Wigner functions are quasidistributions}

Neither \cref{eq:nuclear-wigner-density} nor
\cref{eq:partonic-wigner-density} is required to be pointwise positive.
Negative regions encode interference.  Positivity is imposed on the
appropriate density operator or forward cut correlator before or after the
Wigner transform, not on the Wigner function itself.

\section{Convex geometry and positivity at the correct layer}
\label{sec:positivity}

\subsection{Forward amplitude-level Gram representation}

For a forward correlator with a valid cut or probability interpretation at
fixed kinematics, collect target and parton helicity labels into a composite
index \(\alpha\).  Inserting a complete set of intermediate states gives a
Gram form
\begin{equation}
 \Phi_{\alpha\beta}(x,\kT)
 =\sum_X A_{\alpha X}(x,\kT)
 A^*_{\beta X}(x,\kT),
\label{eq:gram-correlator}
\end{equation}
so that
\begin{equation}
 \Phi=\Phi^\dagger,
 \qquad
 \Phi\succeq0.
\label{eq:forward-psd}
\end{equation}
This density-matrix organization underlies standard spin-1 positivity analyses \cite{Cotogno2017}.  A numerical realization may retain the physical amplitudes \(A\) or use a
block-symmetry-preserving factorization
\begin{equation}
 \Phi=LL^\dagger.
\label{eq:cholesky-correlator}
\end{equation}
Hermiticity and positive semidefiniteness then hold for every parameter
sample without eigenvalue clipping.

Named TMD coefficients are linear coordinates obtained from \(\Phi\).  If
the full matrix variables are retained with linear constraints, the feasible
set is a spectrahedron.  After projecting out unobserved matrix entries, the
allowed named-function region is generally a convex spectrahedral shadow,
not a Cartesian box of independent intervals.

\subsection{Where positivity is not automatic}

The statement \(\Phi\succeq0\) is attached to the operator layer for which
the Gram representation is valid.  A renormalized, rapidity-subtracted, or
finite-order matched TMD is scheme dependent, and a perturbative truncation
need not preserve every pointwise tree-level positivity inequality.
Therefore the architecture distinguishes:
\begin{enumerate}
\item the amplitude-level or low-resolution positive boundary correlator;
\item the exactly transformed correlator under operations known to preserve
positivity;
\item the renormalized/matched/evolved scheme object, for which positivity is
imposed only when a theorem or verified map establishes it, and otherwise is
a diagnostic rather than an unconditional hard constraint.
\end{enumerate}

\begin{requirement}[No clipping]
Negative eigenvalues may not be repaired by replacing them with zero after
independently generated TMDs are assembled.  Such clipping is nonlinear,
changes all correlated projections, and obscures the inconsistent amplitude,
matching map, or truncation that produced the violation.
\end{requirement}

\subsection{Uncertainty samples inside the physical cone}

When positivity applies, uncertainty propagation acts on amplitudes,
Hamiltonian parameters, or a symmetry-adapted factor \(L\), not on independent
named TMD coordinates.  Every ensemble member therefore remains in the same
physical cone while preserving cross-flavor, cross-polarization, and
cross-observable correlations.

\subsection{Information geometry is optional and downstream}

Quantum Fisher information, relative entropy, or other distinguishability
metrics can be useful for experiment design after the physical density
operator, likelihood, and nuisance model are fixed.  They may identify
measurements that distinguish two microscopic mechanisms rather than repeat
an already constrained direction.  They are not substitutes for the
likelihood or for the operator definition.

\section{Matched nucleon-to-deuteron composition}
\label{sec:nuclear-composition}

\subsection{Adopted scale hierarchy}

The first realistic generative implementation uses two matched dynamical
resolutions, in the spirit of controlled nuclear EFT and light-front Hamiltonian matching \cite{Epelbaum2009,Cao2026}:
\begin{enumerate}
\item a quark--gluon light-front model at a hadronic matching scale
\(\Lambda_h\), which generates nucleon states, quark/gluon correlators,
Wilson-line interactions, and matched hadronic operators;
\item a nuclear light-front Hamiltonian or EFT at a lower resolution
\(\Lambda_N\), with explicit \(N\), \(\pi\), \(\Delta\), and any supported
short-distance or non-nucleonic sectors, together with currents and partonic
operator insertions matched from the first layer.
\end{enumerate}
The hierarchy is schematically
\begin{equation}
 H_{\LF}^{qg}(\Lambda_h)
 \xrightarrow{\text{solve}}
 \{|\Psi_N\rangle,\cO_{a/N}^{qg}\}
 \xrightarrow{\cZ_{qg\to\mathrm{had}}}
 \{H_{\LF}^{\mathrm{nuc}}(\Lambda_N),
 \cO_{a}^{\mathrm{nuc}}\}
 \xrightarrow{\text{solve}}
 \{|\Psi_D\rangle,W_{a/D}\}.
\label{eq:hierarchical-nuclear-route}
\end{equation}
This route respects the large separation between internal nucleon dynamics
and deuteron binding while keeping the operator matching, pion content,
short-distance terms, and subtraction rules explicit.

\subsection{Direct quark--gluon deuteron route}

An alternative future calculation solves for a deuteron directly in a
quark--gluon Hilbert space containing six-quark, gluon, sea, cluster, and
hadronic asymptotic sectors.  In that case, the nuclear EFT is a matched
low-resolution reduction of the same full state.  The direct route is more
unified but much more demanding.  It is not required to begin the program,
and results from unrelated nucleon and nuclear microscopic models may not be
combined without an explicit matching interface.

\subsection{Spin-resolved nuclear state}

At the nuclear resolution, write
\begin{equation}
 |D,\Lambda\rangle
 =\Psi_{NN}^{\Lambda}|NN\rangle
 +\Psi_{NN\pi}^{\Lambda}|NN\pi\rangle
 +\Psi_{\Delta\Delta}^{\Lambda}|\Delta\Delta\rangle
 +\Psi_{6q}^{\Lambda}|6q\rangle
 +\cdots.
\label{eq:deuteron-sector-expansion}
\end{equation}
All retained components share one normalization and plus-momentum ledger,
\begin{align}
 1&=Z_{NN}+Z_{NN\pi}+Z_{\Delta\Delta}+Z_{6q}+\cdots,
\label{eq:deuteron-probability-ledger}\\
 1&=\sum_\alpha\int[\dd\Gamma_\alpha]
 \left(\sum_{i\in\alpha}x_i\right)
 |\Psi_\alpha|^2,
\label{eq:deuteron-momentum-ledger}
\end{align}
within the chosen normalization and regulator convention.  Individual sector
probabilities are resolution dependent, but total state normalization,
charge, baryon number, and matched momentum sum rules must close.

The long-distance \(NN\) component must reproduce or improve the successful
S/D phenomenology of high-quality deuteron interactions and must use currents
consistent with the selected Hamiltonian.  Higher sectors require explicit
spin, isospin, color, momentum, and operator content; a scalar probability
alone cannot predict tensor TMDs.

\subsection{Matched nuclear operators}

The QCD partonic operator is represented in the nuclear theory by a matched
expansion
\begin{equation}
 \cO_a^{\mathrm{QCD}}
 \longmapsto
 \cO_a^{\mathrm{nuc}}
 =\sum_{N=p,n}C_{a/N}\otimes\cO_N
 +C_{a/\pi}\otimes\cO_\pi
 +C_{a/NN}\otimes\cO_{NN}
 +C_{a/\Delta}\otimes\cO_\Delta
 +\cdots,
\label{eq:nuclear-operator-matching}
\end{equation}
where the coefficients and operators retain spin, flavor, color-link,
off-forward, and transverse-rank labels.  One-body impulse, exchange-current,
coherent, pion, shadowing, and short-distance terms are controlled reductions
of \cref{eq:nuclear-operator-matching}; they are not independently normalized
curves added after the state is solved.

A pion cloud already represented inside the matched nucleon state and an
explicit \(NN\pi\) nuclear sector generally overlap in resolution \cite{Miller2014}.  The matching must
therefore include a subtraction or replacement rule.  The same requirement
applies to explicit \(qqqg\) content versus an effective gluon input, and to
short-range six-quark operators versus phenomenological SRC response.

\subsection{Full deuteron matrix element}

The deuteron GTMD is evaluated before an incoherent reduction,
\begin{equation}
 W_{a/D}^{[\Gamma,\mathfrak W]}
 =\sum_{\alpha,\beta}
 \langle\Psi_D^\alpha|
 \cO_{a,\alpha\beta}^{\mathrm{nuc},[\Gamma,\mathfrak W]}
 |\Psi_D^\beta\rangle.
\label{eq:full-deuteron-matrix-element}
\end{equation}
Diagonal terms and physically allowed inter-sector interference are retained.
At small \(x\), coherent multiple scattering and shadowing must be computed
from helicity amplitudes consistent with the state and diffractive input
\cite{Frankfurt2012}.
Polarized and tensor shadowing cannot be copied from the unpolarized response.
At intermediate and large \(x\), off-shell, Fermi-motion, SRC, and EMC-like
behavior must arise from the matched joint amplitude and its controlled
separation of scales.

\subsection{Amplitude embedding and completely positive reduction}

Suppose an incoherent active-nucleon channel can be isolated after the full
nuclear amplitude is formed.  The corresponding dilation and Kraus representation follow the standard Stinespring/Kraus construction \cite{Stinespring1955,Kraus1983}.  Introduce an amplitude embedding
\begin{equation}
 V_D:\Hil_N\longrightarrow
 \Hil_D\otimes\Hil_{\mathrm{spec}},
\label{eq:stinespring-embedding}
\end{equation}
where spectator, recoil, partial-wave, and retained channel labels are
explicit.  For a selected channel \(V_D^\dagger V_D\preceq I\); equality holds for a complete isometric embedding.  The induced reduced map is
\begin{equation}
 \cE_D(\rho_N)
 =\Tr_{\mathrm{spec}}
 \left[V_D\rho_NV_D^\dagger\right]
 =\sum_\alpha K_\alpha\rho_NK_\alpha^\dagger.
\label{eq:kraus-from-amplitude}
\end{equation}
It is completely positive by construction.  It is trace preserving only when
all outcomes associated with the embedding are included,
\begin{equation}
 \sum_\alpha K_\alpha^\dagger K_\alpha=I.
\label{eq:kraus-completeness}
\end{equation}
A selected mechanism may instead be trace nonincreasing.

Coherent shadowing, relative phases between Fock sectors, exchange currents,
and any observable sensitive to their interference must be calculated before
the partial trace.  The architecture never forces such physics into an
incoherent Kraus representation.

\section{Symmetry-preserving tensor-network realization}
\label{sec:tensor-network}

\subsection{Network semantics}

A tensor network is used as a sparse representation of amplitudes and
operator contractions.  It is not an assertion that QCD is one-dimensional
or uniformly low entanglement.  An edge label has the form
\begin{equation}
 e=(\nu,R_c,R_f,j^z,\lambda,L_z,\ell_T,\alpha,r,\mathfrak s),
\label{eq:tensor-edge-label}
\end{equation}
where \(\nu\) is a Fock or nuclear sector, \(R_c\) and \(R_f\) are color and
flavor representations, \(j^z,\lambda,L_z\) are angular labels, \(\ell_T\) is the transverse tensor rank,
\(\alpha\) is a basis multiplicity, \(r\) is the regulator/truncation label,
and \(\mathfrak s\) is any operator scheme metadata carried by the edge.

Nodes are physical maps: Hamiltonian vertices, invariant couplers,
wave-function coefficients, currents, Wilson-line interactions, matched
operator insertions, nuclear sector couplers, or controlled partial traces.
Every tensor index must map to a physical degree of freedom or a declared
numerical basis label.

\subsection{Shared-parameter propagation}

A parameter belongs to one Hamiltonian term, current, matching coefficient,
state tensor, or discrepancy operator.  It propagates automatically to every
observable contraction that contains that object.  Automatic differentiation
or adjoint methods give
\begin{equation}
 J_{ia}=\frac{\partial\cO_i}{\partial\theta_a},
 \qquad
 \operatorname{Cov}(\cO)
 \simeq J\operatorname{Cov}(\theta)J^T
\label{eq:autodiff-covariance}
\end{equation}
for the parameter component of the uncertainty, while preserving exact
symmetry blocks.

\subsection{Gauge and symmetry blocks}

Block sparsity is generated from intertwiners and conservation laws, not from
post hoc pruning.  Color tensors are contracted to gauge-invariant external
states; spin-1 spherical tensors are converted to Cartesian TMD channels by a
fixed invertible adapter; transverse-rank blocks carry the correct
\(SO(2)\) weight; and forbidden blocks are absent rather than penalized.

\subsection{Convergence obligation}

Low bond dimension or low factorization rank can suppress precisely the OAM,
color, sea, coherent, or long-range correlations sought by the model.
Network rank is therefore one truncation axis among several.  Results must be
reported versus bond/factorization rank together with basis, regulator, Fock,
and quadrature convergence.  The network is a contraction strategy, not a
replacement for \(H_{\LF}\).

\section{Provenance 2-complex and no-double-counting audit}
\label{sec:provenance}

\subsection{The typed composition graph is primary}

The first implementation remains a directed typed graph.  Vertices are
states, operator spaces, matched descriptions, or observables.  Edges are
physical or descriptive maps with source, target, scheme, replacement, and
normalization metadata.  A path records the ancestry of a final contribution.
The graph rejects a composition if the endpoint types do not match.

\subsection{Cellular construction}

To make alternative derivations and subtraction identities auditable, extend
the graph to a two-dimensional cellular complex \(K\):
\begin{itemize}
\item \(0\)-cells are typed state/operator/observable objects;
\item oriented \(1\)-cells are declared transitions, matching maps,
reductions, or composition steps;
\item \(2\)-cells are attached only when physics declares two paths
comparable, equivalent up to a matching coefficient, or related by an
explicit subtraction/replacement identity.
\end{itemize}
The cellular chain groups satisfy
\begin{equation}
 C_2(K)\xrightarrow{\partial_2}C_1(K)
 \xrightarrow{\partial_1}C_0(K),
 \qquad
 \partial_1\partial_2=0.
\label{eq:provenance-chain-complex}
\end{equation}
For a two-cell whose boundary is the difference of two source-to-target
paths, \(\partial_2\sigma=p_1-p_2\), the cell records the declared physical
relation between those paths.

\subsection{How the audit detects double counting}

Examples include:
\begin{enumerate}
\item a pion cloud already contained in the matched nucleon state versus an
explicit nuclear \(NN\pi\) contribution;
\item a gluon distribution encoded in an effective nucleon input versus an
explicit \(qqqg\) sector;
\item a screened eikonal phase versus a separately parameterized direct
imaginary amplitude intended to represent the same rescattering;
\item an integrated-out Fock sector versus the counterterm or effective
operator that replaces it.
\end{enumerate}
If two paths represent the same contribution, the corresponding two-cell must
specify whether one replaces the other, whether a subtraction is required, or
whether they are mutually exclusive ensemble alternatives.  The composition
engine flags a final sum that traverses both paths without the declared rule.

\subsection{Interpretation of homology}

The group
\begin{equation}
 H_1(K)=\ker\partial_1/\operatorname{im}\partial_2
\label{eq:provenance-homology}
\end{equation}
identifies path cycles not removed by the declared equivalence and subtraction
cells.  Such a cycle is an audit signal: it may represent an unresolved
alternative, a missing unitarity relation, an unclosed replacement rule, or a
genuinely distinct interference mechanism.  It is not asserted that a
physical amplitude depends only on a homology class.

\begin{principle}[Limited topological claim]
The provenance complex organizes composition identities.  It does not turn
QCD dynamics, Wilson-line phases, or nuclear amplitudes into topological
invariants.  Higher categorical or homological structures are introduced only
when a concrete equivalence, obstruction, or computational need supplies the
cells and their physical meaning.
\end{principle}

\section{Regulator, matching, and convergence towers}
\label{sec:truncation}

\subsection{Directed trial spaces}

Let a truncation label be
\begin{equation}
 r=(\Lambda_{\mathrm{UV}},\Lambda_{\mathrm{IR}},N_{\max},K,
 \mathcal F,\chi,\ldots),
\label{eq:truncation-label}
\end{equation}
where \(\mathcal F\) is the retained Fock set and \(\chi\) may denote tensor
network rank or another compression parameter.  The calculations form a
directed family \(\{\Hil_r\}_{r\in\mathfrak R}\).  When two basis choices
are literally nested, an inclusion
\begin{equation}
 I_{r\to r'}:\Hil_r\hookrightarrow\Hil_{r'}
\end{equation}
is available.  When the regulator, basis, or degrees of freedom change, no
canonical inclusion is assumed; instead the model supplies comparison or
prolongation maps
\begin{equation}
 C_{r\to r'}:\Hil_r\longrightarrow\Hil_{r'}
\label{eq:state-comparison-map}
\end{equation}
with a documented approximation meaning.

Bare eigenvectors at different regulators are not required to converge in a
common Hilbert-space norm.  The primary convergence objects are matched
masses, current matrix elements, GTMD moments, and process observables.

\subsection{Inverse system of effective operators}

Each level has an effective operator algebra \(\Alg_r\) and a renormalized
mass operator
\begin{equation}
 \cM_r^2=\cM_{0,r}^2+\sum_a c_a(r)\cO_{a,r}.
\label{eq:effective-mass-operator}
\end{equation}
Coarse-graining and matching maps
\begin{equation}
 R_{r'\to r}:\Alg_{r'}\longrightarrow\Alg_r
\label{eq:operator-inverse-system}
\end{equation}
form an inverse system only to the accuracy at which the EFT or regulator
matching is constructed.  For a calibration observable,
\begin{equation}
 R_{r'\to r}[\cO_{r'}]
 =\cO_r+\delta_{r'r},
 \qquad
 \delta_{r'r}\xrightarrow[r,r'\to\mathrm{continuum}]{}0
\label{eq:projective-observable-consistency}
\end{equation}
according to a fitted or power-counted truncation law.

\subsection{Required convergence axes}

A predictive claim must test stability under:
\begin{enumerate}
\item longitudinal and transverse basis enlargement;
\item UV and IR regulator variation with counterterm refitting;
\item addition of the next physically required Fock sector;
\item restoration tests for Poincar\'e and current constraints;
\item Wilson-line/eikonal order;
\item nuclear many-body and current order;
\item tensor-network rank or other compression;
\item quadrature, interpolation, transform range, and grid refinement;
\item perturbative matching order and scale variation.
\end{enumerate}
Convergence in one direction cannot hide nonconvergence in another.

\subsection{Observable-level convergence metric}

For a vector of matched observables \(\bm\cO_r\), define
\begin{equation}
 d(r,r')=
 \left[
 (\bm\cO_r-\bm\cO_{r'})^T
 \Sigma_{\mathrm{ref}}^{-1}
 (\bm\cO_r-\bm\cO_{r'})
 \right]^{1/2},
\label{eq:observable-convergence-distance}
\end{equation}
where \(\Sigma_{\mathrm{ref}}\) specifies the physically relevant scales and
correlations.  A convergence statement reports \(d(r,r')\), componentwise
residuals, and the extrapolation model.  Smoothness of a single curve at one
truncation is not evidence of convergence.

\section{Inference, uncertainty, and falsifiability}
\label{sec:inference}

\subsection{Shared microscopic and matching parameters}

Let
\begin{equation}
 \theta=
 (\theta_{qg},\theta_{\mathrm{nuc}},\theta_{\mathrm{current}},
 \theta_{\mathrm{match}},\theta_{\mathrm{Wilson}})
\label{eq:shared-parameter-vector}
\end{equation}
contain quark/gluon masses and couplings, low-energy constants, counterterms,
current coefficients, matching coefficients, and a restricted set of
rescattering parameters.  Each component is attached to a specific physical
operator or controlled discrepancy term.  No latent parameter is introduced
only because one named TMD is poorly described.

For calibration data \(D_{\mathrm{cal}}\),
\begin{equation}
 p(\theta,\delta\mid D_{\mathrm{cal}})
 \propto
 p(D_{\mathrm{cal}}\mid\theta,\delta)
 p(\theta)p(\delta),
\label{eq:microscopic-posterior}
\end{equation}
where \(\delta\) contains explicitly separated EFT, regulator, matching,
nuclear, Wilson-line, and numerical discrepancy components.

Suitable calibration information includes masses, charges, magnetic and
axial moments, electromagnetic and energy--momentum form factors, selected
PDFs and low moments, matched lattice matrix elements, deuteron binding and
asymptotic S/D information, elastic moments, a restricted subset of
\(b_1\), SIDIS, and T-odd data, and any process for which the factorization
and operator matching are controlled.

\subsection{Withheld predictions}

Before final optimization, entire observable families are divided into
calibration and prediction sets.  The posterior predictive distribution is
\begin{equation}
 p(D_{\mathrm{hold}}\mid D_{\mathrm{cal}})
 =\int\dd\theta\,\dd\delta\,
 p(D_{\mathrm{hold}}\mid\theta,\delta)
 p(\theta,\delta\mid D_{\mathrm{cal}}).
\label{eq:posterior-predictive}
\end{equation}
Examples of meaningful tests include:
\begin{itemize}
\item calibrating a shared rescattering kernel with selected Sivers data and
predicting Boer--Mulders, tensor T-odd, or gluon color-sector observables;
\item calibrating masses, currents, and selected collinear data and predicting
nonzero-transfer GTMD moments and TMD/GPD marginal closure;
\item calibrating the nucleon model and long-distance deuteron properties,
then predicting tensor-polarized deuteron observables without a new
normalization per function.
\end{itemize}

A failed withheld family must revise the state, Hamiltonian, current,
matching, factorization assumption, or truncation/discrepancy model.  Adding a
coefficient directly to the failed TMD returns the construction to a
patchwork.

\subsection{Correlated uncertainty decomposition}

For observables \(\cO_i\),
\begin{equation}
 \operatorname{Cov}(\cO_i,\cO_j)
 =\left\langle
 (\cO_i-\langle\cO_i\rangle)
 (\cO_j-\langle\cO_j\rangle)
 \right\rangle_{\theta,\delta}.
\label{eq:observable-covariance}
\end{equation}
The ensemble preserves cross-flavor, quark--gluon, proton--neutron,
TMD--GPD--current, nucleon--deuteron, and cross-process correlations.
The error budget reports separately:
\begin{enumerate}
\item statistical/calibration uncertainty;
\item quark--gluon Hamiltonian/EFT truncation;
\item Fock-sector and basis truncation;
\item UV/IR regulator and renormalization uncertainty;
\item Wilson-line/eikonal truncation;
\item quark--gluon-to-hadronic matching uncertainty;
\item perturbative QCD/TMD matching and scale variation;
\item Collins--Soper evolution and large-\(b_{\mathrm{TMD}}\) uncertainty;
\item nuclear interaction, current, and many-body truncation;
\item numerical discretization, transform, and solver error.
\end{enumerate}
Positivity is a constraint on an applicable state space, not an uncertainty
estimate.

\section{Software realization and formal interfaces}
\label{sec:software}

\subsection{Core objects}

\begin{longtable}{L{0.27\textwidth}L{0.65\textwidth}}
\toprule
Object & Responsibility\\
\midrule
\endhead
\texttt{ConventionContract} & Metric, light-front normalization, Fourier
signs, momentum fractions, state normalization, \(\bD\)/\(\bM\)/\(\RT\)
types, and unit conventions.\\
\texttt{RegulatedSectorSpace} & Momentum simplex, basis, Fock grade,
constituent representations, regulator, and physical projectors.\\
\texttt{GaugeConstraint} & Gauss-law/color-singlet projection, color basis,
and gauge-covariance tests.\\
\texttt{Intertwiner} & Sparse invariant tensor with domain, codomain,
multiplicity, symmetry action, and equivariance tests.\\
\texttt{LFMassOperator} & Renormalized block mass-squared operator
\(2P^+P^- -\bm P_T^2\), counterterms, sector-changing vertices, and solver
adapter.\\
\texttt{FockState} & Normalized amplitudes, momentum fiber, phase convention,
sector probabilities, member identity, and truncation metadata.\\
\texttt{BareWilsonPath} & Endpoint fibers, path geometry, representation,
orientation, concatenation, inverse, and thin-homotopy/reparametrization
metadata.\\
\texttt{BareHolonomy} & Gauge-covariant parallel transport satisfying the
bare path functor laws.\\
\texttt{RenormalizedLink\allowbreak Operator} & Bare transport plus cusp, rapidity,
soft-subtraction, scale, and scheme data.\\
\texttt{GTMDOperator} & Momentum-fiber-changing quark/gluon insertion and
common Fock-overlap evaluator at nonzero \(\DT\).\\
\texttt{MatchingMap} & Quark--gluon-to-hadronic EFT matching, Hamiltonian-to-
QCD operator matching, OPE coefficients, and truncation remainder.\\
\texttt{ReductionMap} & Typed forward limit, Wigner transform, GPD reduction,
local moment, and rank-aware TMD Fourier transform.\\
\texttt{PositiveBoundary\allowbreak Correlator} & Amplitude/Gram or symmetry-adapted
\(LL^\dagger\) representation at layers where positivity is valid.\\
\texttt{NuclearLFModel} & Matched nuclear Hamiltonian, \(NN\), \(NN\pi\),
\(\Delta\Delta\), and supported short-distance sectors with consistent
currents.\\
\texttt{NuclearAmplitude\allowbreak Embedding} & Full coherent amplitude map and any
subsequent Stinespring/Kraus reduction with trace ledger.\\
\texttt{TensorContractionGraph} & Block-sparse symmetry-preserving state,
operator, and observable contractions with automatic differentiation.\\
\texttt{TruncationTower} & Directed trial spaces, inverse operator matching,
convergence fits, and discrepancy estimates.\\
\texttt{Provenance2Complex} & Typed graph, declared path-equivalence
2-cells, subtraction/replacement rules, and observable ancestry.\\
\texttt{ObservableKernel} & Process hard factors, fragmentation, color
weights, factorization status, and \(W+Y\) matching.\\
\texttt{PosteriorEnsemble} & Shared microscopic member identity through all
states, operators, reductions, and observables.\\
\bottomrule
\end{longtable}

\subsection{Required property-based tests}

\begin{enumerate}
\item With the declared \(v^\pm=(v^0\pm v^3)/\sqrt2\) convention, the mass
operator is exactly \(2P^+P^- -\bm P_T^2\).
\item \texttt{WignerImpactCoordinate}, \texttt{TMDSeparation}, and
\texttt{NuclearCoordinate} cannot be passed to one another without an explicit
transform or adapter.
\item Physical sector states satisfy color singlet/Gauss constraints,
permutation statistics, momentum support, and total-quantum-number closure.
\item Every intertwiner commutes with the represented exact symmetry action;
controlled breaking operators fail the symmetric test only by their declared
amount.
\item The regulated mass operator is Hermitian, and Poincar\'e/current defects
are tracked versus the truncation tower.
\item Bare holonomy respects path concatenation, endpoint typing, gauge
covariance, and inversion.
\item Renormalized link operators refuse composition when cusp, rapidity,
soft, scale, or scheme labels are incompatible.
\item Time reversal maps future and past link classes with the correct
T-even/T-odd behavior where factorization establishes the relation.
\item Matched GTMD reduction squares commute within separately reported
matching, truncation, and numerical tolerances.
\item The \(\DT\)-Fourier transform uses \(\bD\); the \(\kT\)-Fourier
transform uses \(\bM\); and the nuclear \(\DNT\)-Fourier transform uses \(\RT\).  Cross-use is a type error.
\item Forward boundary correlators remain positive without clipping at layers
where the Gram theorem applies; Wigner functions are exempt from pointwise
positivity.
\item Spin-1 rotations reproduce the \(0\oplus1\oplus2\) transformation law
and the fixed adapter to \(U,L,T,LL,LT,TT\).
\item Disabling either amplitude in an OAM/helicity interference removes the
corresponding transverse harmonic.
\item The nuclear CP map is generated from an amplitude embedding and partial
trace; coherent terms remain outside the trace until physically justified.
\item Total nuclear normalization, charge, baryon number, and plus momentum
close after all declared sectors and matching/subtraction terms are included.
\item The provenance complex detects deliberately duplicated pion, gluon,
sea, phase, and short-distance mechanisms.
\item Process assembly refuses missing hard/color factors, unsupported gauge
links, or an undeclared factorization status.
\end{enumerate}

\subsection{Role of quantum simulation}

PennyLane or another quantum simulator is justified only after a finite
physical basis, Hamiltonian matrix, observables, and encoding are specified.
A quantum circuit must reproduce exact diagonalization in the same truncated
space before being accepted for variational eigensolving, differentiation, or
correlation diagnostics.  A generic entangling circuit that reproduces only
Clebsch--Gordan coupling or adds untyped latent parameters does not add
predictive physics.

\section{Validation ladder}
\label{sec:validation}

Validation proceeds in increasing dynamical depth.  Passing a later numerical
fit does not excuse failure at an earlier structural layer.

\begin{enumerate}
\item \textbf{Convention and type layer:} light-front normalization, Fourier
signs, momentum support, unit consistency, distinct transverse-coordinate
types, and fail-closed scheme composition.
\item \textbf{Gauge and algebraic layer:} color singlets, Gauss constraints,
Hermiticity, exact represented symmetries, controlled breaking, parity,
time reversal with link reversal, permutation symmetry, and spin-1 projection
closure.
\item \textbf{State layer:} normalized eigenstates, mass spectrum, regulator
and basis studies, rotational multiplets, and stable matched observables as
Fock content is enlarged.
\item \textbf{Current layer:} charges, magnetic and axial moments,
energy--momentum moments, Ward identities, angular conditions, and currents
consistent with the Hamiltonian.
\item \textbf{Common-parent layer:} nonzero-\(\DT\) GTMDs and matched
TMD/GPD/current/Wigner reductions from the same amplitudes.
\item \textbf{Nucleon layer:} withheld flavor-resolved PDFs, TMD moments,
lattice matrix elements, form factors, OAM-sensitive relations, and
process-link reversal.
\item \textbf{Nuclear layer:} deuteron binding, S/D observables, elastic
moments, \(b_1\), tagged structure, tensor moments, plus-momentum closure,
and controlled proton/neutron and no-additional-sector limits.
\item \textbf{Process layer:} SIDIS, Drell--Yan, inclusive DIS, and
factorization-qualified gluon-sensitive predictions with explicit hard,
fragmentation, color, evolution, and \(W+Y\) kernels.
\item \textbf{Posterior predictive layer:} successful prediction of at least
one nontrivial observable family excluded from calibration.
\end{enumerate}

A failed observable is traced backward through the provenance graph and
matched reduction squares.  The permitted responses are to revise a physical
interaction, state sector, current, matching coefficient, factorization
assumption, or discrepancy model.  A new TMD-specific normalization is not a
permitted default repair.

\section{Staged research program}
\label{sec:stages}

\subsection{Stage 0: convention, type, and scale-separation contract}

Before changing central physics:
\begin{enumerate}
\item implement the light-front convention of
\cref{eq:lf-components,eq:lf-mass-operator};
\item create distinct software types for \(\DT,\bD,\kT,\bM,\DNT,\pT,\RT,\qT\);
\item define \(\Amp\), \(\Dens\), \(\Match\), \(\Red\), and \(\Proc\)
interfaces and forbid cross-layer implicit composition;
\item make Wilson path, representation, cusp, regulator, soft, scale, and
scheme data part of operator identity;
\item adopt the hierarchical quark--gluon-to-nuclear matching route of
\cref{eq:hierarchical-nuclear-route};
\item define the first truncation labels and observable-level convergence
metrics.
\end{enumerate}

\subsection{Stage A: migrate the accepted boundary without changing its physics}

\begin{enumerate}
\item encode the current accepted quark and gluon correlator parents as
symmetry-adapted positive boundary tensors where positivity applies;
\item encode all existing registry projectors as typed intertwiners and
reduction maps;
\item encode present composition and replacement rules in the provenance
graph and initial two-cells;
\item reproduce the existing accepted numerical parents, marginals, closure
residuals, and benchmark observables to regression tolerance;
\item demonstrate that the new architecture detects deliberately injected
scheme, coordinate, rank, and double-counting errors.
\end{enumerate}
This stage validates the architecture but makes no new microscopic claim.

\subsection{Stage B: microscopic nucleon pilot}

Select one primary regulated light-front route and solve at least
\begin{equation}
 |qqq\rangle\oplus|qqqg\rangle\oplus|qqqq\bar q\rangle
\end{equation}
at multiple basis, regulator, and Fock truncations.  Compute a nonzero-
\(\DT\) T-even subset of quark and gluon GTMDs, their matched TMD/GPD/current
reductions, and straight-link OAM moments.  The pilot is successful only if
one shared parameter set describes or predicts more than one observable
family and the matched reduction squares close.

\subsection{Stage C: dynamical Wilson-line pilot}

Add a factorization-appropriate eikonal interaction derived from the same
gauge sector.  Test zero-rescattering limits, future/past reversal,
flavor/operator differences, gluon \(f/d\) contractions, one-gluon matching,
higher-order stability, and the dependence on rapidity/IR regulation.  No
universal fitted imaginary phase is permitted.

\subsection{Stage D: matched spin-1 nuclear model}

Match the nucleon dynamics and partonic operators onto a nuclear light-front
Hamiltonian/EFT.  Reproduce binding, asymptotic S/D information, static
moments, elastic form factors, and consistent one-/two-body current tests.
Add \(NN\pi\) and any other supported sectors with one normalization and
explicit subtraction rules.  Reproduce controlled \(b_1\) limits before
claiming predictions for the complete tensor TMD sector.

\subsection{Stage E: common QCD scheme, evolution, and process kernels}

Match every retained operator into one declared QCD TMD/GTMD scheme.  Evolve
nonsinglet, singlet, quark--gluon, and all transverse-rank sectors with the
appropriate quark or gluon Collins--Soper kernel and threshold treatment.
Add process hard factors, fragmentation where required, color weights,
factorization-status labels, and process-specific \(W+Y\) matching.

\subsection{Stage F: inference and withheld validation}

Calibrate the restricted shared parameter set, propagate all truncation and
matching axes, and reserve at least one nontrivial nucleon or deuteron family
for posterior predictive validation.  This is the first stage at which the
phrase ``genuinely predictive next-level model'' is auditable.

\section{Risks, failure modes, and explicit nonclaims}
\label{sec:risks}

\begin{longtable}{L{0.22\textwidth}L{0.34\textwidth}L{0.36\textwidth}}
\toprule
Risk & Failure mode & Required safeguard\\
\midrule
\endhead
Decorative category theory & Calling every operation a morphism in one
category despite incompatible composition laws. & Maintain separate amplitude,
CP, matching, reduction, and process layers; validate only typed mixed
squares.\\
Decorative topology & Calling link classes, cycles, or TMDs topological
invariants without dynamical derivation. & Restrict topology to bare path
organization and the provenance 2-complex; retain connection, curvature,
Hamiltonian, and matching dependence.\\
Light-front convention error & Using \(P^+P^-\) as the mass squared with the
\(1/\sqrt2\) convention. & Freeze the convention contract and test
\(\cM^2=2P^+P^- -\bm P_T^2\).\\
Transverse-variable collision & Identifying \(\bD\) with \(\bM\), or a
nuclear coordinate with either. & Distinct types and transform-specific
interfaces; no implicit casts.\\
Gauge-category overstatement & Treating an open Wilson line as a gauge-
invariant endomorphism or the renormalized TMD as a bare path functor. & Use
endpoint fibers for bare transport and a decorated renormalized operator type.\\
Naive marginal equality & Equating a literal \(\kT\) integral of a
soft-subtracted TMD to a PDF/current. & Use renormalized operator matching and
the small-\(\bM\) OPE in commuting squares.\\
Symmetry overconstraint & Treating color gauge symmetry, isospin, spin, and
\(L_z\) as identical exact global factors. & Separate gauge constraints,
approximate internal symmetry, and light-front Poincar\'e structure; include
breaking operators.\\
Tensor-network bias & Low rank suppresses long-range, OAM, sea, or color
correlations. & Publish network-rank convergence with basis and Fock
convergence.\\
Positivity misuse & Enforcing pointwise positivity on Wigner
quasidistributions or every finite-order scheme object. & Apply Gram/PSD
constraints only at layers where the theorem is valid; propagate scheme
caveats explicitly.\\
Channel misuse & Forcing coherent shadowing or sector interference into an
incoherent CP map. & Form the full amplitude first; trace only after the
observable no longer resolves the coherence.\\
Nuclear patchwork at a deeper level & Combining a microscopic nucleon with an
unrelated phenomenological nuclear response. & Use the matched hierarchy and
consistent currents/operators; compare alternative nuclear theories as
controlled model variants.\\
Double counting across scales & Counting pion, gluon, sea, or short-distance
physics both in constituent structure and explicit sectors. & Require
matching/subtraction two-cells and composition-engine enforcement.\\
False state convergence & Comparing bare wave functions or Fock
probabilities across unrelated regulators as observables. & Compare matched
masses and matrix elements; treat decompositions as scheme diagnostics.\\
Parameter migration & Hiding one free coefficient per TMD inside sophisticated
notation. & Attach parameters only to interactions, currents, matching
coefficients, or explicit discrepancy operators.\\
Factorization overreach & Using a TMD formula in a color-entangled or
unproven channel. & Carry factorization status in every process object and
refuse unsupported universal links or hard factors.\\
\bottomrule
\end{longtable}

The proposal does not claim:
\begin{itemize}
\item an exact continuum solution of QCD;
\item that a tensor network supplies missing interactions;
\item that Wilson-line dependence is topological in generic nonzero
curvature;
\item that all renormalized TMD schemes preserve pointwise positivity;
\item that every nuclear mechanism is an incoherent quantum channel;
\item that one direct quark--gluon deuteron Hamiltonian must be solved before
a matched predictive hierarchy is possible;
\item that successful evolution of the present phenomenological boundary by
itself changes the model into a microscopic prediction.
\end{itemize}

\section{Recommended architecture and acceptance criteria}
\label{sec:decision}

\subsection{Recommended foundation}

The adopted foundation is a
\begin{quote}
\emph{gauge-covariant, graded light-front tensor architecture with typed
amplitude, density, matching, reduction, and process layers; matched
quark--gluon-to-nuclear scale separation; common GTMD operators; bare path
transport and decorated renormalized Wilson operators; phase-space and
positive-state geometry at their correct layers; and a directed
regulator/truncation tower.}
\end{quote}

Its indispensable elements are:
\begin{enumerate}
\item the convention and type contract, including distinct \(\bD\) and
\(\bM\);
\item regulated physical light-front sectors and the correctly normalized
mass-squared operator;
\item gauge constraints and representation-theoretic intertwiners;
\item common nonzero-transfer GTMD operators and matched reductions;
\item bare Wilson transport plus decorated renormalized link operators;
\item the hierarchical nucleon-to-nuclear matching route with consistent
currents and subtraction rules;
\item amplitude-derived positive forward correlators where valid;
\item full coherent nuclear amplitudes with CP reductions only after justified
partial traces;
\item a symmetry-preserving tensor-network contraction layer;
\item the provenance graph/2-complex and directed convergence tower;
\item shared-parameter inference and genuinely withheld predictions.
\end{enumerate}

\subsection{Acceptance criteria}

The next-level model may be called predictive only when all of the following
conditions are met:
\begin{enumerate}
\item One declared matched hierarchy, or one direct quark--gluon Hamiltonian,
generates the central quark, antiquark, and gluon correlators.  Presently
fitted named-TMD curves are calibration or comparison data, not indispensable
central inputs.
\item Nucleon sea and gluon observables arise from retained Fock sectors or
matched effective operators with demonstrated regulator, basis, and omitted-
sector control.
\item Proton/neutron flavor differences, helicity, transversity, OAM, and
spin--orbit correlations arise from the state and declared breaking
operators.
\item T-even GTMDs and their TMD/GPD/current/form-factor/Wigner reductions
close through matched diagrams from the same amplitudes.
\item T-odd phases arise from an explicit Wilson-line/rescattering
calculation with correct process reversal and gluon \(f/d\) color structure.
\item The spin-1 nuclear state contains normalized supported sectors with
explicit tensor polarization and physically allowed interference.
\item Nuclear interactions, currents, and partonic operators are matched
consistently; impulse, exchange, coherent, and non-nucleonic labels denote
controlled reductions rather than independent normalizations.
\item One declared QCD scheme and rank-aware evolution route carries all
retained components to multiple \(Q\) values without erasing provenance or
closure.
\item A shared microscopic posterior propagates correlated nucleon and
deuteron uncertainties, with truncation and matching axes reported
separately.
\item At least one nontrivial observable family excluded from calibration is
successfully predicted within the declared uncertainty model.
\item Removing every presently phenomenological named-TMD input does not make
the central microscopic calculation undefined.
\item All convention, gauge, symmetry, positivity-applicability, marginal,
sum-rule, regulator, truncation, numerical, factorization, and
no-double-counting gates pass.
\end{enumerate}

\section{Conclusion}

The next advance is not a larger atlas of separately parameterized functions.
It is a change in the mathematical and dynamical source of the functions.
The corrected architecture makes that change precise.  It fixes the
light-front normalization, separates the two transverse impact variables,
distinguishes gauge constraints from approximate internal symmetries,
separates amplitude maps from completely positive reductions and matching
maps, and replaces naive marginal equalities by renormalized commuting
squares.  It treats bare Wilson transport and renormalized TMD operators as
related but different objects, imposes positivity only where the operator
admits a Gram representation, and derives nuclear CP maps from amplitude
embeddings rather than calling amplitudes themselves completely positive.

Most importantly, it adopts a realizable matched hierarchy: quark--gluon
nucleon dynamics generate common partonic amplitudes and operators; these are
matched onto a spin-resolved nuclear light-front theory with consistent
currents, sectors, and subtraction rules; deuteron GTMDs and observables are
then matrix elements of the resulting state.  A direct quark--gluon deuteron
solution remains a deeper future limit, not a prerequisite that prevents
progress.

If implemented faithfully, the architecture converts many current validation
tests into structural consequences and makes the remaining failures
scientifically diagnostic.  The value of the algebra and geometry is not that
they solve QCD by terminology, but that they prevent flavor, spin, OAM,
gauge-link class, nuclear composition, renormalization, positivity, and
uncertainty from drifting into mutually inconsistent models.

\appendix

\section{Compact mathematical dictionary}
\label{app:dictionary}

\begin{tabularx}{\textwidth}{L{0.31\textwidth}Y}
\toprule
Physical concept & Mathematical/computational object\\
\midrule
Light-front convention & \texttt{ConventionContract};
\(v^\pm=(v^0\pm v^3)/\sqrt2\),
\(\cM^2=2P^+P^- -\bm P_T^2\)\\
Fock sectors & Graded direct sum of regulated physical Hilbert spaces\\
Momentum support & Longitudinal simplex and transverse constraint manifold\\
Color gauge invariance & Associated-bundle representations plus Gauss/color-
singlet projector\\
Flavor, target spin, parton helicity, OAM & Representation labels and
irreducible blocks\\
Allowed interaction & Equivariant intertwiner in the amplitude layer\\
Spin-1 vector/tensor structure &
\(\End(\Rep_1)\cong\Rep_0\oplus\Rep_1\oplus\Rep_2\)\\
Sparse contraction & Symmetry-preserving block tensor network\\
Bare gauge link & Parallel transport between color-bundle endpoint fibers\\
Staple/process class & Decorated morphism in an admissible path groupoid\\
Renormalized TMD link & Bare path plus cusp, rapidity, soft, scale, and scheme
data\\
GTMD & Momentum-fiber-changing gauge-covariant operator matrix element\\
Partonic Wigner impact coordinate & \(\bD\), Fourier conjugate to \(\DT\)\\
TMD separation & \(\bM\), Fourier conjugate to \(\kT\)\\
Nuclear Wigner coordinate & \(\RT\), Fourier conjugate to nuclear transfer \(\DNT\), with average momentum \(\pT\) retained\\
TMD/GPD/current/Wigner relations & Typed matched reduction maps and commuting
squares\\
Transverse OAM organizer & Moment map on \((\bD,\kT)\in T^*\mathbb R^2\)\\
Forward positivity & Gram/PSD cone at the applicable amplitude-level operator\\
Nuclear impulse channel & Amplitude embedding followed, where valid, by a CP
partial-trace map\\
Nuclear scale separation & Matched quark--gluon-to-hadronic light-front
hierarchy\\
Double-counting audit & Typed provenance graph with declared equivalence and
subtraction 2-cells\\
Truncation sequence & Directed trial spaces plus inverse effective-operator
matching system\\
Predictive uncertainty & Pushforward of shared posterior and separated
truncation/matching discrepancies\\
\bottomrule
\end{tabularx}

\section{Minimal notation ledger}
\label{app:notation}

\begin{longtable}{L{0.20\textwidth}L{0.70\textwidth}}
\toprule
Symbol & Meaning\\
\midrule
\endhead
\(\kT\) & Average intrinsic transverse momentum of the active parton in the
GTMD/TMD correlator.\\
\(\DT\) & Transverse momentum transfer between external hadron states.\\
\(\bD\) & Partonic impact coordinate Fourier conjugate to \(\DT\).\\
\(\bM\) & Bilocal TMD transverse separation Fourier conjugate to \(\kT\).\\
\(\DNT\) & Off-forward transverse transfer in the nuclear spectral density; it is related to \(\DT\) by the declared recoil map.\\
\(\pT\) & Average internal active-nucleon transverse momentum retained in the nuclear Wigner transform.\\
\(\RT\) & Nuclear impact coordinate Fourier conjugate to \(\DNT\).\\
\(\qT\) & Measured process-level transverse momentum.\\
\(\mathfrak W\) & Fully decorated Wilson/operator type, including path,
representation, regulators, soft subtraction, scales, and scheme.\\
\(\cZ_{qg\to\mathrm{had}}\) & Matching from quark--gluon degrees of freedom
and operators to the nuclear hadronic EFT.\\
\(\cZ_{\LF\to\mathrm{QCD}}\) & Matching from the regulated light-front
operator to the declared renormalized QCD/GTMD/TMD scheme.\\
\(\Amp,\Dens,\Match\)\newline\(\Red,\Proc\) & Distinct typed composition layers for
amplitudes, density reductions, matching, marginals, and process assembly.\\
\bottomrule
\end{longtable}

\section{Adoption questions}
\label{app:questions}

\begin{enumerate}
\item Can every tensor index be mapped to a physical degree of freedom or a
controlled numerical basis label?
\item Does every parameter belong to a Hamiltonian term, current, matched
operator, Wilson interaction, or explicit discrepancy model?
\item Are \(\DT\), \(\bD\), \(\kT\), \(\bM\), \(\DNT\), \(\pT\), and \(\RT\) separate types throughout the code?
\item Does every state use the declared light-front normalization and the
correct invariant mass operator?
\item Are local gauge constraints separated from approximate isospin and
light-front rotational structure?
\item Are bare path transport and renormalized link operators distinct?
\item Can at least two matched observable reductions be computed from the
same nonzero-transfer correlator?
\item Is positivity structural at the layer where it is valid and absent as a
false requirement where it is not?
\item Does every CP nuclear map descend from an explicit amplitude embedding
and justified trace?
\item Are all pion, sea, gluon, phase, SRC, and non-nucleonic overlaps covered
by matching/subtraction rules?
\item Is every approximation embedded in a directed convergence or matching
sequence?
\item Does the architecture enable a genuinely withheld cross-channel
prediction?
\end{enumerate}
A negative answer identifies a concrete physics, matching, or implementation
problem.  It must not be hidden by renaming the object.

\small
\begin{thebibliography}{99}

\bibitem{Dirac1949}
P.~A.~M. Dirac,
``Forms of relativistic dynamics,''
Rev. Mod. Phys. \textbf{21}, 392 (1949).

\bibitem{Brodsky1998}
S.~J. Brodsky, H.-C. Pauli, and S.~S. Pinsky,
``Quantum chromodynamics and other field theories on the light cone,''
Phys. Rept. \textbf{301}, 299 (1998),
arXiv:hep-ph/9705477.

\bibitem{Vary2009}
J.~P. Vary et al.,
``Hamiltonian light-front field theory in a basis function approach,''
Phys. Rev. C \textbf{81}, 035205 (2010),
arXiv:0905.1411.

\bibitem{Meissner2009}
S.~Mei{\ss}ner, A.~Metz, and M.~Schlegel,
``Generalized parton correlation functions for a spin-\(\tfrac12\) hadron,''
JHEP \textbf{08}, 056 (2009),
arXiv:0906.5323.

\bibitem{Belitsky2004}
A.~V. Belitsky, X.~Ji, and F.~Yuan,
``Quark imaging in the proton via quantum phase-space distributions,''
Phys. Rev. D \textbf{69}, 074014 (2004),
arXiv:hep-ph/0307383.

\bibitem{LorcePasquini2011}
C.~Lorc\'e and B.~Pasquini,
``Quark Wigner distributions and orbital angular momentum,''
Phys. Rev. D \textbf{84}, 014015 (2011),
arXiv:1106.0139.

\bibitem{Lorce2012}
C.~Lorc\'e,
``Wilson lines and orbital angular momentum,''
Phys. Lett. B \textbf{719}, 185 (2013),
arXiv:1210.2581.

\bibitem{BacchettaMulders2000}
A.~Bacchetta and P.~J. Mulders,
``Deep inelastic leptoproduction of spin-one hadrons,''
Phys. Rev. D \textbf{62}, 114004 (2000),
arXiv:hep-ph/0007120.

\bibitem{KumanoSong2021}
S.~Kumano and Q.-T. Song,
``Transverse-momentum-dependent parton distribution functions for spin-1
hadrons,''
Phys. Rev. D \textbf{103}, 014025 (2021),
arXiv:2011.08583.

\bibitem{Boer2016}
D.~Boer, S.~Cotogno, T.~van Daal, P.~J. Mulders, A.~Signori, and Y.-J. Zhou,
``Gluon and Wilson loop TMDs for hadrons of spin \(\leq1\),''
JHEP \textbf{10}, 013 (2016),
arXiv:1607.01654.

\bibitem{Cotogno2017}
S.~Cotogno, T.~van Daal, and P.~J. Mulders,
``Positivity bounds on gluon TMDs for hadrons of spin \(\leq1\),''
JHEP \textbf{11}, 185 (2017),
arXiv:1709.07827.

\bibitem{Collins2011}
J.~Collins,
\emph{Foundations of Perturbative QCD},
Cambridge University Press (2011).

\bibitem{Echevarria2012}
M.~G. Echevarria, A.~Idilbi, and I.~Scimemi,
``Factorization theorem for Drell--Yan at low \(q_T\) and
transverse-momentum distributions on the light cone,''
JHEP \textbf{07}, 002 (2012),
arXiv:1111.4996.

\bibitem{Neff2016}
T.~Neff and H.~Feldmeier,
``The Wigner function and short-range correlations in the deuteron,''
arXiv:1610.04066.

\bibitem{Epelbaum2009}
E.~Epelbaum, H.-W. Hammer, and U.-G. Mei{\ss}ner,
``Modern theory of nuclear forces,''
Rev. Mod. Phys. \textbf{81}, 1773 (2009),
arXiv:0811.1338.

\bibitem{Miller2014}
G.~A. Miller,
``Pionic and hidden-color, six-quark contributions to the deuteron
\(b_1\) structure function,''
Phys. Rev. C \textbf{89}, 045203 (2014),
arXiv:1311.4561.

\bibitem{Frankfurt2012}
L.~Frankfurt, V.~Guzey, and M.~Strikman,
``Leading twist nuclear shadowing phenomena in hard processes with nuclei,''
Phys. Rept. \textbf{512}, 255 (2012),
arXiv:1106.2091.

\bibitem{Stinespring1955}
W.~F. Stinespring,
``Positive functions on \(C^*\)-algebras,''
Proc. Amer. Math. Soc. \textbf{6}, 211 (1955).

\bibitem{Kraus1983}
K.~Kraus,
\emph{States, Effects, and Operations},
Lecture Notes in Physics 190, Springer (1983).

\bibitem{MacLane1998}
S.~Mac Lane,
\emph{Categories for the Working Mathematician}, 2nd ed.,
Springer (1998).

\bibitem{Selinger2007}
P.~Selinger,
``Dagger compact closed categories and completely positive maps,''
Electron. Notes Theor. Comput. Sci. \textbf{170}, 139 (2007).

\bibitem{Cao2026}
X.~Cao et al.,
``Non-perturbative flavor asymmetry in the nucleon and deuteron: the
light-front Hamiltonian effective field theory approach,''
arXiv:2601.13567 (2026).

\end{thebibliography}

\end{document}
